%%%%%%%%%%%%%%%%%%%%%%%%%%%%%%%%%%%%% FORMAT
%%%%%%%%%%%%%%%%%%%%%%%%%%% footer and header
% \usepackage{lastpage}
% \usepackage{fancyhdr}
% \usepackage{xparse}
% \pagestyle{fancy}
%%%%%%%%%%%%%%%%%%%%%%%%%%% Headers and footers 

% \fancyhead{} % clear all header fields
% % \renewcommand{\headrulewidth}{0pt} % no line in header area
% \fancyfoot{} % clear all footer fields
% % \fancyfoot[LE,RO]{\thepage}           % page number in "outer" position of footer line


% \rfoot{\thepage\ of \pageref{LastPage}}
% % Define the footer separation line
% \renewcommand{\footrulewidth}{0.4pt}
% \usepackage[headheight=65pt,tmargin=65pt,headsep=5pt, footskip=20pt]{geometry}
% \geometry{a4paper, margin=2cm}

% \usepackage[numbers]{natbib}
 % \usepackage[comma]{natbib}
% \usepackage[natbib=true]{biblatex}

%\let\labelindent\relax
%\let\proof\relax
%\let\endproof\relax

%%%%%%%%%%%% Floating-base dynamics

% numActuatedJoints
\newcommand{\ajDim}{n}
\newcommand{\fbDim}{6}
\newcommand{\jsDim}{(\ajDim+\fbDim)}
\newcommand{\actuationMatrix}{\text{B}}
% \newcommand{\numActuatedJoints}{n}
% \newcommand{\fbDim}{6}
% \newcommand{\jsDim}{(\numActuatedJoints+\fbDim)}

%%%%%%%%%%%%%%%%%%% Units
% \usepackage{xspace}
\newcommand{\unitPos}{~m\xspace}

\newcommand{\unitLM}{~kg \cdot m/s\xspace}
\newcommand{\unitMass}{~kg\xspace}
\newcommand{\unitHertz}{~N/m$^{3/2}$\xspace}

\newcommand{\unitCm}{~cm\xspace}

\newcommand{\unitAM}{~kg \cdot m^2/s\xspace}
\newcommand{\unitPosJS}{~rad\xspace}
\newcommand{\unitVelTS}{~m/s\xspace}
\newcommand{\unitVelJS}{~rad/s\xspace}
\newcommand{\unitAccTS}{~m/s$^2$\xspace}
\newcommand{\unitAccVel}{~rad/s$^2$\xspace}
\newcommand{\unitForce}{~N\xspace}
\newcommand{\unitImpulse}{(N$\inner{}{}$s)\xspace}
\newcommand{\unitMS}{~ms\xspace}
\newcommand{\unitFrequency}{~Hz\xspace}



%%%%%%%%%%%%%%%%%%%%%%%%%%% Table
% \usepackage{multirow}

%%%%%%%%%%%%%%%%%%%%%%%%%%% Font colour 

% \usepackage[utf8]{inputenc}
% \usepackage[english]{babel}
%\usepackage[dvipsnames]{xcolor}
% \usepackage{listings}
% \usepackage{color}
% %%%%%%%%%%%%%%%%%%%%%%%%% Tikz packages:
% \usepackage{tikz}
% \usepackage{tikz-layers}
% \usepackage{array}
% \usepackage{subfigure}
% \newcolumntype{C}[1]{>{\centering\arraybackslash}p{#1}}
%
% the positioning-plus library (is on TeX.sx),
% the node-families library (is on TeX.sx),
% the fit library (loaded by positioning-plus),
% the backgrounds library to draw stuff behind other stuff,
% the calc library for some funny coordinate calculations and the
% shapes.geometric library for the ellipse shape

% \usetikzlibrary{tikz-layers}
\usetikzlibrary{shapes.geometric,automata,positioning,arrows,shadows,,backgrounds,calc}

\definecolor{frameColor}{RGB}{128,128,128}
\definecolor{initialStateColor}{RGB}{10,100,10}
\definecolor{impactStateColor}{RGB}{0,255,0}
\definecolor{admittanceStateColor}{RGB}{0,0,255}

\definecolor{resetStateColor}{RGB}{10,100,10}

\definecolor{stateColor}{RGB}{51,204,204}

\definecolor{controllerColor}{RGB}{10,10,255}
\definecolor{plannerColor}{RGB}{10,150,150}
\definecolor{estimatorColor}{RGB}{10,200,10}
\definecolor{robotColor}{RGB}{255,0,0}
\definecolor{modelColor}{RGB}{180,100,10}

%%% Enumeration
% \usepackage{enumitem, kantlipsum}
% \usepackage{enumitem}

% \usepackage{kantlipsum}

%%%%%%%%%%%%%%%%%%%%%%%%%%% Adjust font size
% \usepackage{setspace} 

%%%%%%%%%%%%%%%%%%%%%%%%%%%%%%%%%%%%%% MATH:



%%%%%%%%%%%%%%%%%%%%%%%%% AMS packages:
% \usepackage{mathtools}
% \usepackage{amsmath, amsthm, amsfonts}
% \usepackage{amssymb}
% \usepackage{graphicx}
% \usepackage[font=normal,labelfont=bf]{caption}
% \captionsetup{font=footnotesize}
% \captionsetup[table]{format=plain,labelformat=simple,labelsep=period}%
% \usepackage{subcaption}
% \usepackage{url}
% \usepackage{accents} %for symbols below a character, e.g. \underaccent{\bar}
% \usepackage[colorlinks=true,hyperindex=true,bookmarksnumbered,bookmarks=true]{hyperref} % Turn on "bookmarks" and "hyperindex" to generate outlines. 

%%%%%%%%%%%%%%%%%%%%%%%%%% Self-defined commands: 


\newcommand{\setDef}[2]{
  \defeq \{#1 : #2 \}
}

\newcommand{\quickEq}[2]{
  \begin{equation}
    \label{#1}
    {#2}
  \end{equation}
}

\newcommand{\quickAd}[1]{
  \begin{aligned}
    {#1}
  \end{aligned}
}

\newcommand{\de}[1]{
  \frac{d}{dt}(#1)
}
%%%%%%%%%%%%%%%%%%%%%%%%%% notes with colors:
\newcommand{\az}[1]{\textcolor{red}{[AZ:#1]}}
%%%%%%%%%%%%%%%%%%%%%%%%%% Frames:
\def\frameCamera{$\cal R_{\rm camera}$ }
\def\frameRobot{$\cal R_{\rm robot}$ }
\def\frameBox{$\cal R_{\rm box}$ }
\def\frameSurface{$\cal R_{\rm s_i}$ }
%%%%%%%%%%%%%%%%%%%%%%%%%% Units:
\newcommand{\unitHz}{~Hz}
\newcommand{\unitMs}{~ms}

%%%%%%%%%%%%%%%%%%%%%%%%%% LQR commands: 
\newcommand{\costFunctional}{J}
\newcommand{\stateWeight}{Q}

\newcommand{\inputWeight}{R}

\newcommand{\discountF}{\gamma}

% Set
% \newcommand{\set}{\sigma}
\newcommand{\setSymbol}{\mathcal{X}}
\newcommand{\set}[1]{\setSymbol_{#1}}

%%%%%%%%%%%%%%%%%%%%%%%%%% Optimal Control

\newcommand{\costate}{\lambda}%
\newcommand{\hamiltonian}{H}%

\newcommand{\initialTime}{t_0}%
\newcommand{\finalTime}{t_f}%

% #1 denotes the input weight: R; #2 denotes the state
\newcommand{\lqrControl}[2]{
  - \inverse{(#1 + \quadraticObj{B}{P_{n+1}})}\transpose{B}P_{n+1}A#2
}%

\newcommand{\manifold}{\mathcal{S}}%


%%%%%%%%%%%%%%%%%%%%%%%%%% Impact dynamics:
% \usepackage{pifont}% http://ctan.org/pkg/pifont
\newcommand{\cmark}{\ding{51}}%
\newcommand{\xmark}{\ding{55}}%

% Contact Vel:
\newcommand{\contactVel}{\bs{v}}
\newcommand{\nContactVel}{v}


% notation of the impacting body 
\newcommand{\ib}{1}

% notation of the object 
\newcommand{\ob}{2}

% virtual spring length
\newcommand{\vSpringL}{x}
\newcommand{\vSpringLd}{\dot{x}}
% spring coefficient 
\newcommand{\vSpringC}{k}

% The potential energy stored in the virtual spring
\newcommand{\pEnergy}{{E}}
% #1 denotes the normal impulse
\newcommand{\pEnergyIntOne}[1]{\Phi_1(#1)}
% #1 denotes the normal impulse
\newcommand{\pEnergyIntTwo}[1]{\Phi_2(#1)}

\newcommand{\pEnergyd}{\dot{E}}


% Impulse at the end of the impact process
\newcommand{\impulseEnd}{\impulse_r}

% Derivative w.r.t. the normal impulse
\newcommand{\dni}[1]{\frac{d #1}{d \nImpulse} }
\newcommand{\dzi}[1]{\frac{d #1}{d \zImpulse} }

% \newcommand{\di}[1]{{\prime{#1}}}


% Inverse inertia matrix 
\newcommand{\iim}{W}

% Energetic coefficient of restitution
\newcommand{\eCoefR}{e_{\text{r}}}

\newcommand{\optimal}[1]{{#1}^{*}}

\newcommand{\postImpact}[1]{{#1}^{+}}
\newcommand{\preImpact}[1]{#1^{-}}
\newcommand{\coefR}{c_{\text{r}}}
% Coefficient of friction
\newcommand{\coefF}{\mu}

% The impulse when the restitution phase ends
\newcommand{\impulseRE}{\impulse_{r}}

% tangential direction unit vector
\newcommand{\tangentialU}{\hat{\bs{t}}}

% normal direction unit vector
\newcommand{\normalU}{\hat{\bs{n}}}

% Z direction of the inertial frame
\newcommand{\nz}{\text{\bf n}}

\newcommand{\fConeNum}{N_\coefF}

% The tangential impulse
\newcommand{\tImpulse}{\impulseScalar_{t}}
\newcommand{\tImpulsed}{\dot{\impulseScalar}_{t}}
% The tangential force
\newcommand{\tForce}{\forceScalar_{t}}

% The tangential linear Vel
\newcommand{\tVel}{\linearVelScalar_{t}}
% The normal linear Vel
\newcommand{\nVel}{\linearVelScalar_{n}}

% The normal impulse
\newcommand{\nImpulse}{{\impulseScalar_{n}}}
\newcommand{\nImpulsed}{\dot{\impulseScalar}_{n}}
\newcommand{\zImpulse}{{\impulseScalar_{z}}}
% The  unit vector of a coordinate system:
\newcommand{\basisVec}[1]{\widehat{\bs{\bs{#1}}}}

% The normal force
\newcommand{\nForce}{\forceScalar_{n}}
% The tangential impulse

\newcommand{\tpImpulse}{\tp{\impulse}}


% Denote the tangential plane:
\newcommand{\tp}[1]{{#1}_{\perp}}


% The normal impulse when the compression phase ends
\newcommand{\nImpulseCE}{\impulseScalar_{nc}}
\newcommand{\nImpulseCEini}{\tilde{\impulseScalar}_{nc}}

% The normal impulse when the restitution phase ends
\newcommand{\nImpulseRE}{\impulseScalar_{nr}}

% The smallest value of normal impulse at which v_t = 0.
\newcommand{\nImpulseS}{\impulseScalar_{ns}}

% The derivative of impulse w.r.t. normal impulse before I_ns
\newcommand{\iDniB}{\bs{\delta}}

% The derivative of impulse w.r.t. normal impulse after I_ns
\newcommand{\iDniA}{\bs{\sigma}}


% The potential energy when compression phase ends
\newcommand{\pEnergyC}{\pEnergy_{c}}

% The potential energy when v_t = 0 first holds
\newcommand{\pEnergyT}{\pEnergy_{s}}



%%%%%%%%%%%%%%%%%%%%%%%%%% MPC commands: 
\definecolor{csrStateColor}{RGB}{255,255,153}
\definecolor{scrStateColor}{RGB}{153,204,0}
\definecolor{crStateColor}{RGB}{153,204,155}

\newcommand{\footTask}{p}

% Foot Step Location:
\newcommand{\fL}{p}
% timeStep
\newcommand{\timeStep}{N}
% footStep
\newcommand{\footStep}{M}


% footStep
\newcommand{\stateVar}{\bs{x}}

% #1 variable with different steps, #2 step number
\newcommand{\step}[2]{
  {#1^{#2}}
}

%%%%%%%%%%%%%%%%%%%%%%%%%% Eigenvalue:
\newcommand{\eigV}{
\sigma
}

\newcommand{\eigVec}{
\bs{e}
}



%%%%%%%%%%%%%%%%%%%%%%%%%% Rigid body motion


\newcommand{\Vee}[1]{\left(#1\right)^\vee}


%%%%%%%%% Represent the derivative operator
\newcommand{\derivative}[1]{\frac{d}{dt}(#1)}
%%% #1:function, #2:var1, #3:var2
\newcommand{\pdv}[3]{
  \frac{\partial^2 #1}{\partial #2 \partial #3}
}
%%% #1:function, #2:var1
\newcommand{\pd}[2]{
  \frac{\partial #1}{\partial #2 }
}

\newcommand{\guassD}{\mathcal{D}}

\newcommand{\body}[1]{#1^{b}}
\newcommand{\spatial}[1]{#1^{s}}

\newcommand{\jointPosition}{\bs{\theta}}
\newcommand{\jointPositiond}{\dot{\jointPosition}}
\newcommand{\jointPositiondd}{\ddot{\jointPosition}}
\newcommand{\fbJointPosition}{\bs{\mathcal{X}}}

% The frame defined by Christian Ott: 
\newcommand{\coFrame}{\mathcal{C}}

\newcommand{\hybridFrame}{H}
\newcommand{\twistVel}{\bs{V}}

% The instantaneous inertial frame (whose origin locates on the ground and shares the COM's X-Y coordinates).
\newcommand{\iiFrame}{I}

\newcommand{\fbFrame}{B}
\newcommand{\inertialFrame}{\text{O}}
\newcommand{\inertialWrench}{{\wrench_{\inertialFrame}}}
\newcommand{\inertialForce}{{\force_{\inertialFrame}}}
\newcommand{\inertialTorque}{{\torque_{\inertialFrame}}}
\newcommand{\comFrame}{\com}

% Gravity force
\newcommand{\gravityandcoriolis}{\mathbf{N}}
\newcommand{\gForce}{\bs{g}}
\newcommand{\gValue}{\text{9.81}}
\newcommand{\gForceDef}{\vectorThree{0}{0}{-1}g}

\newcommand{\cmmFrame}{G}
% \newcommand{\cmmFrame}{G}

\newcommand{\fbVel}{
  \bs{V}_{\fbFrame}
}
% \newcommand{\fbVelDef}{
%   \spatialVel{\inertialFrame}{\fbFrame}
% }
% \newcommand{\fbVelDefDef}{
%   \spatialVelDef{\inertialFrame}{\fbFrame}
% }

\newcommand{\cmmVel}{
  \textcolor{bodyVelColor}{
    \bs{V}_{\cmmFrame}
  }
}

\newcommand{\cmmVelDef}{
  \bodyVel{\inertialFrame}{\cmmFrame}
}
\newcommand{\cmmVelDefDef}{
  \inverse{\cmmInertia} \cmmMomenta
}
%%%%%%%%%%%% Exponential Map given unit skew matrix and a real number
%%% #1:input matrix
\newcommand{\epMapDef}[1]{
\identityMatrix + #1  + \frac{#1^2}{2!}  + \frac{#1^3}{3!} + \ldots
}
%%% #1:unit skew matrix #2 realnumber 
\newcommand{\unitSkewExpDef}[2]{
\identityMatrix + #1sin(#2) + #1^2(1-cos(#2))
}

%%%%%%%%%%% -------------------------------------- Jacobian:
%%%% #1 reference frame #2 body frame #3 joint angle
\newcommand{\transformConfig}[3]{
  \textcolor{geometricColorVariable}{
    \transform{#1}{#2}(#3)
  }}
\newcommand{\transformConfigDef}[3]{
  \textcolor{geometricColorVariable}{
    \epMap{\twist{#1}{#2}#3} \transform{#1}{#2}(0)
  }}

\newcommand{\transformdConfigDef}[3]{
  \textcolor{geometricColorVariable}{
    \twist{#1}{#2} \dot{#3}\epMap{\twist{#1}{#2}#3} \transform{#1}{#2}(0)
  }}
\newcommand{\transformdConfigDefTwo}[3]{
  \textcolor{geometricColorVariable}{
    \epMap{\twist{#1}{#2}#3}\twist{#1}{#2} \dot{#3} \transform{#1}{#2}(0)
  }}

\newcommand{\transformddConfigDef}[3]{
  \textcolor{geometricColorVariable}{
    % \twist \ddot{#3}\epMap{\twist#3} \transform{#1}{#2}(0)
    % +\twist \dot{#3}\twist \dot{#3}\epMap{\twist#3} \transform{#1}{#2}(0)
    (\twist{#1}{#2} \ddot{#3}%\epMap{\twist#3} \transform{#1}{#2}(0)
    +\twist{#1}{#2} \twist{#1}{#2} \dot{#3}^2)\epMap{\twist{#1}{#2}#3} \transform{#1}{#2}(0)
  }}
\newcommand{\transformddConfigDefTwo}[3]{
  \textcolor{geometricColorVariable}{
    % \twist \ddot{#3}\epMap{\twist#3} \transform{#1}{#2}(0)
    % +\twist \dot{#3}\twist \dot{#3}\epMap{\twist#3} \transform{#1}{#2}(0)
    \epMap{\twist{#1}{#2}#3} (\twist{#1}{#2} \ddot{#3}%\epMap{\twist#3} \transform{#1}{#2}(0)
    +\twist{#1}{#2} \twist{#1}{#2} \dot{#3}^2)\transform{#1}{#2}(0)
  }}

% #1:reference frame, #:2 end-effector frame
\newcommand{\poeDef}[2]{
  \textcolor{geometricColorVariable}{
    \twistEp{#1}{1}{1}\twistEp{#1}{2}{2}\ldots\twistEp{#1}{#2}{#2}\transformConfig{#1}{#2}{\zeroVector}
  }}
\newcommand{\invPoeDef}[2]{
  \textcolor{geometricColorVariable}{
    \transformInvConfig{#1}{#2}{\zeroVector}\invTwistEp{#1}{#2}{#2}\ldots\invTwistEp{#1}{2}{2}\invTwistEp{#1}{1}{1}
  }}

%%%% #1 reference frame #2 body frame #3 joint angle
\newcommand{\transformInvConfig}[3]{
  \textcolor{geometricColorVariable}{
    \transformInv{#1}{#2}(#3)
  }}
\newcommand{\transformdConfig}[3]{
  \textcolor{geometricColorVariable}{
    \transformd{#1}{#2}(#3)
  }}
\newcommand{\transformddConfig}[3]{
 \textcolor{geometricColorVariable}{
   \transformdd{#1}{#2}(#3)
 }}
\newcommand{\transformInvConfigDef}[3]{
  \textcolor{geometricColorVariable}{
  \transformInv{#1}{#2}(0) \epMap{-\twist{#1}{#2}#3}
  }}



%%% Lie bracket
\newcommand{\twoCases}[4]{
\begin{dcases}
  {#1} & {#2} \\
  {#3} & {#4}
\end{dcases}
}
% \newcommand{\up}[1]{\skewMatrix{#1}}

\newcommand{\liebra}[2]{
  [#1,#2]
}
\newcommand{\liebraResult}[2]{
  \Vee{\liebraDefDef{#1}{#2}}
}
%%%% #1:linear velocity #2: angular velocity #3:linear velocity #4: angular velocity 
\newcommand{\seLiebraResult}[4]{
(\cross{#2}{#3} - \cross{#4}{#1}, \skewMatrix{#2}#4 - \skewMatrix{#4}#2 )
}



\newcommand{\liebraDef}[2]{
  \skewMatrixTwo{\cross{#1}{#2}} 
}

\newcommand{\liebraDefDef}[2]{
  #1#2 - #2#1
}

%%%%%%%%%%% ----------------------------------------Scew
%%%%%%%%%%%% Transform defined by a screw
%%% #1:axis, #2:theta, #3:point, #4:pitch(finite)
\newcommand{\screwTransform}[4]{
\matrixTwo{
  \epMap{\skewMatrix{#1}#2}
}{
  (\identityMatrix - \epMap{\skewMatrix{#1}#2})#3
  + #4 #1 #2
}{
\zeroVector
}{
1
}
}

%%%%%%%%%%%% Wrench defined by a screw
%%% #1:magnitude, #2:point, #3:direction, #4:pitch(finite)
%%% #1: reference frame #2: screw name
\newcommand{\geoWrenchDef}[2]{
  \screwM_{#2} \vectorTwo{
    \angularVel_{#2}
  }{
    -\cross{\angularVel_{#2}}{\translation{\origin{#1}}{#2}} + \pitch_{#2}\angularVel_{#2}
  }
}

%%%%%% When the pitch is infinite:
%%% #1: screw name
\newcommand{\geoWrenchDefTwo}[1]{
  \screwM_{#1} \vectorTwo{
    \zeroVector
  }{
    \angularVel_{#1}
  }
}

\newcommand{\distance}{d}
%%% #1:screw 1, #2:screw 2 
\newcommand{\screwAngle}[2]{
  \alpha_{#1#2}
}
\newcommand{\screwAngleDef}[2]{
atan2(\inner{(\cross{\angularVel_{#1}}{\angularVel_{#2}})}{\unitVec_{#1#2}}, \inner{\angularVel_{#1}}{\angularVel_{#2}} ).
}
\newcommand{\unitVec}{\bs{n}}

% Calculates the reciprocal product
\newcommand{\recip}[2]{
  #1 \odot{}   #2
}

\newcommand{\recipScrewDef}[2]{
  \screwM_{#1}\screwM_{#2}\left(
    (\pitch_{#1} + \pitch_{#2})\cos\screwAngle{#1}{#2} - \distance\sin\screwAngle{1}{2}
  \right)
}

%%%%%%%%%%%% Pitch:
\newcommand{\pitch}{
  h
}



% Twist2Screw:  #1 linear vel #2 angular vel
\newcommand{\pitchTwistDef}[2]{
  \frac{\innerP{#2}{#1}}{\twoNorm{#2}}
}

% Wrench2Screw: #1 force #2 moment
\newcommand{\pitchWrenchDef}[2]{
  \frac{\innerP{#1}{#2}}{\twoNorm{#1}}
}


\newcommand{\twoNorm}[1]{
  {{\left\lVert#1\right\rVert}^2}
}

\newcommand{\norm}[1]{
  {\left\lVert#1\right\rVert}
}
  \newcommand{\abs}[1]{\left\vert#1\right\vert}
  \newcommand{\cross}[2]{ #1 \times #2 }
\newcommand{\quantity}{\bs{d}}
\newcommand{\quantityd}{\dot{\bs{d}}}
\newcommand{\quantitydd}{\ddot{\bs{d}}}
\newcommand{\state}{\bs{x}}

%%%%%%%%%%%% Axis:
\newcommand{\axis}{
  l
}

%%% #1:point, #2:angularVel direction
\newcommand{\screwAxisDef}[2]{
  \{ #1 + \lambda#2 : \lambda \in \RRv{} \}
}

% When angularVel is Not zero: #1:linearVel, #2:angularVel, #3:lambda (scalar) 
\newcommand{\axisTwistDefThree}[3]{
\frac{\cross{#2}{#1}}{\twoNorm{#2}} + \axisTwistDefTwo{#1}{#3}
}

% When angularVel is zero: #1:linearVel, #2:lambda (scalar) 
\newcommand{\axisTwistDefTwo}[2]{
   #1 #2
}

% When force is NOT zero: #1:force, #2:torque, #3:lambda (scalar) 
\newcommand{\axisWrenchDefThree}[3]{
\frac{\cross{#1}{#2}}{\twoNorm{#1}} + \axisWrenchDefTwo{#2}{#3}
}

% When force is zero: #1 torque #2: lambda (scalar) 
\newcommand{\axisWrenchDefTwo}[2]{
   #1 #2
}

%%%%%%%%%%%% Magnitude:
\newcommand{\screwM}{
  M
}

% Wrench2Screw: #1: force,  if force is nonzero. Otherwise #1 is the torque
% Twist2Screw: #1: angularVel,  if angularVel is nonzero. Otherwise #1 is the linearVel
\newcommand{\screwMDef}[1]{
  \norm{#1}
}
%%%%%%%%%%%% Twist:

% #1 axis #2 theta 
\newcommand{\twistE}[2]{
  e^{\skewMatrix{#1}#2}
}


% Computes the iner product
\newcommand{\innerP}[2]{
  \transpose{#1}#2
}

% Computes the outer product
\newcommand{\outerP}[2]{
  #1\transpose{#2}
}



% Calculates the reciprocal product between two twists
% #1:v_1,  #2:w_1,  #3:v_2, #4:w_2.
\newcommand{\recipVelDef}[4]{
  \innerP{#1}{#4} + \innerP{#3}{#2}
}

% Calculates the reciprocal product between two wrenches
% #1:force,  #2:torque,  #3:force, #4:torque. 
\newcommand{\recipWrenchDef}[4]{
  \innerP{#1}{#4} + \innerP{#3}{#2}
}

% Screw
\newcommand{\screw}{S}
%%% #1 screw notation
\newcommand{\screwDef}[1]{
(\screwM_{#1}, \pitch_{#1}, \point{#1}, \angularVel_{#1})
}


% Detail the wrench along a screw
\newcommand{\wrenchAS}[2]{S}

% Detail the twist about a screw
%%% #1:axis, #2:\point, #3 pitch
\newcommand{\linearTwistAS}[3]{
-\cross{#1}{#2} + #3#1
}

\newcommand{\linearTwistASTwo}[3]{
  \cross{#2}{#1} + #3#1
}

% Detail the twist about a screw, when the pitch is zero: 
%%% #1:axis, #2:\point, 
\newcommand{\linearTwistASRevolute}[2]{
-\cross{#1}{#2}
}
\newcommand{\linearTwistASRevoluteTwo}[2]{
{#2}\cross{#1}
}



%%%%%%%%%%% Composite rigid body (CRB) algorithm:
% \usepackage{algorithmicx}
% \usepackage{algorithm}% http://ctan.org/pkg/algorithms
% \usepackage{algpseudocode}% http://ctan.org/pkg/algorithmicx
%%% predecessor:
\newcommand{\pre}[1]{
  p(#1)
}
%%% successor:
\newcommand{\suc}[1]{
  s(#1)
}

%%% CRB vel

\newcommand{\crbVel}{
  {^{crb}\bs{V}}
}
%%% CRB inertia
\newcommand{\crbGInertia}{
  {^{crb}I}
}

%%% System inertia
\newcommand{\systemInertia}{
  \textcolor{bodyVelColor}{
    I
  }
}


%%%%%%%%%%% 

\newcommand{\averageVel}{
  \bar{\systemVel}
}
\newcommand{\relativeVel}{{\systemVel'}}

\newcommand{\dof}{n}
\newcommand{\inertiaMatrix}{\text{M}}

\newcommand{\equivalentMass}{{\inertiaMatrix_{eq}}}
\newcommand{\osInertia}{{\Lambda}}

% #1:reference frame #2:contact index
\newcommand{\equivalentMassDef}[2]{
  \inverse{\left(
      \spatialJacobian{#1}{#2} \inverse{\inertiaMatrix}_{#2}
     \transpose{{\spatialJacobian{#1}{#2}}}
   \right)}
}

\newcommand{\inertiaMatrixd}{\dot{\inertiaMatrix}}

%%% #1:row index, #2:column index
\newcommand{\inertiaMatrixComponentDef}[2]{
\sum^\dof_{l = max(#1,#2)}{
  \transpose{
    \twistC{S}{#1}(t_0)
  }\transpose{
    \twistTransform{#1}{l}
  }
  \spatialGInertia{S}{\com_l}(t_0)
  \twistTransform{#2}{l}\twistC{S}{#2}(t_0)
}
}

\newcommand{\inertiaMatrixComponentDefTwo}[2]{
\sum^\dof_{l = max(#1,#2)}{
  \transpose{
    \twistC{S}{#1}(t_0)
  }\transpose{
    \adg{l}{#1}
  }
  \spatialGInertia{S}{\com_l}(t_0)
  \adg{l}{#2}
  \twistC{S}{#2}(t_0)
}
}

%%%%%%%%%%%%%%%%%% inertiaMatrix partial derivative:
%%% #1:row index, #2:column index, #3 theta index
\newcommand{\inertiaMatrixdComponentDef}[3]{
  \sum^\dof_{l = max(#1,#2)}{\left(
      \transpose{
        \liebra{\adg{#3}{#1}\twistC{S}{#1}}{\twistC{S}{#3}}
      }
      \transpose{
        \adg{l}{#3}
      }
      \spatialGInertia{S}{\com_l}(t_0)
      \adg{l}{#2}
      \twistC{S}{#2}(t_0)
      +
      \transpose{
        \twistC{S}{#1}(t_0)
      }\transpose{
        \adg{l}{#1}
      }
      \spatialGInertia{S}{\com_l}(t_0)
      \adg{l}{#3}\liebra{\adg{#3}{#2}\twistC{S}{#2}}{\twistC{S}{#3}}.
      % \adg{l}{#2}
      % \twistC{S}{#2}(t_0)
    \right)
  }
}



\newcommand{\cMatrix}{C(\jointPosition, \jointPositiond)}
\newcommand{\nVector}{N}
\newcommand{\rpTransform}{\mathcal{T}}
\newcommand{\rpTransformd}[1]{\mathcal{#1}}

\newcommand{\coordinateTransformed}[2]{
\transpose{{\inverse{#2}}} #1 \inverse{#2}
}



\newcommand{\cristoffelSymbol}[3]{\Gamma_{#1#2#3}}

%%% #1 Reference frame: #2 notation
\newcommand{\twist}[2]{
  \textcolor{spatialVelColor}{  
    \skewMatrix{\xi}_{#1#2}
  }
}

%%% #1: infinitesimal translation generator #2:so{3} 
\newcommand{\twistDef}[2]{
  \matrixTwo{
    #2
  }{
    #1
  }{\zeroVector
  }{0
  }
}
%%%%%%%%%%%%%%% The exponential map of a twist:
%% #1 linear part #2 angularPart   #3 theta
\newcommand{\twistEpDef}[3]{
  \matrixTwo{
    \epMap{\skewMatrix{#2} #3}
  }{
    (\identityMatrix - \epMap{\skewMatrix{#2} #3})
    (\cross{#2}{#1})
    +
    \projectionDef{#2}#1#3
  }{
    \zeroVector
  }{
    1
  }
}
%% #1: reference frame, #2 twist notation #3 theta notation
\newcommand{\twistEp}[3]{
\epMap{\twist{#1}{#2} \theta_{#3}}
}
%% #1: reference frame, #2 twist notation #3 theta notation
\newcommand{\invTwistEp}[3]{
\epMap{-\twist{#1}{#2} \theta_{#3}}
}


%%%%%%%%%%%%%%%%%%%%%% Dynamics
%%% #1:index, #2:last index, #3 component
\newcommand{\agg}[3]{
\sum^{#2}_{#1=1}{#3}
}

%%% #1:index, #2:last index, #3 component
\newcommand{\union}[3]{
\bigcup^{#2}_{#1=1}{#3}
}


% %%%%%%%%%%% #1:row index, #2:column index
% \newcommand{\cMatrixComponentDef}[2]{
%   \agg{k}{\dof}{\cMatrixComponent{#1}{#2}{k}\dot{\jointPosition}_k} 
% }

%%%%%%%%%%% #1:row index, #2:column index, #3 index (between 1 and dof)
\newcommand{\cMatrixComponent}[3]{
  \Gamma_{#1#2#3}
}
%%%%%%%%%%% #1:row index, #2:column index, #3 index (between 1 and dof)
\newcommand{\cMatrixComponentDef}[3]{
  \frac{1}{2}\left(
    \pd{\inertiaMatrix_{#1#2}}{\jointPosition_{#3}}
    +\pd{\inertiaMatrix_{#1#3}}{\jointPosition_{#2}}
    -\pd{\inertiaMatrix_{#3#2}}{\jointPosition_{#1}}
  \right)
}



%%% The twist coordinate
% \newcommand{\twistC}{\xi}
% #1:reference frame, #2:notation of the twist 
% \newcommand{\twistC}{\xi}

% #1:reference frame, #2:notation of the twist 
\newcommand{\twistC}[2]{
  \textcolor{spatialVelColor}{
    \xi_{#1#2}
  }
}

% #1:reference frame, #2:notation of the twist 
\newcommand{\twistCd}[2]{
  \textcolor{spatialVelColor}{
    \dot{\xi}_{#1#2}
  }
}

% #1:reference frame, #2:notation of the twist 
\newcommand{\spatialTwist}[2]{
  \textcolor{spatialVelColor}{
    \hat{\xi}'_{#1#2}
  }}

% #1:reference frame, #2:notation of the twist 
\newcommand{\spatialTwistC}[2]{
  \textcolor{spatialVelColor}{
    \xi'_{#1#2}
  }}

% #1:reference frame, #2:notation of the twist 
\newcommand{\spatialTwistCd}[2]{
  \textcolor{spatialVelColor}{
    \dot{\xi}'_{#1#2}
  }}

% #1:reference frame, #2:notation of the twist 
\newcommand{\bodyTwistC}[2]{
  \textcolor{bodyVelColor}{
    \xi^{\dagger}_{#1#2}
  }}
% #1:reference frame, #2:index of the twist 
\newcommand{\bodyTwistCd}[2]{
  \textcolor{bodyVelColor}{
    \dot{\xi}^{\dagger}_{#1#2}
  }}

\newcommand{\bodyTwist}[2]{
  \textcolor{bodyVelColor}{
    \hat{\xi}^{\dagger}_{#1#2}
  }}

% #1:reference frame, #2:index  of the twist, #3:index of the starting link
%%% t_0 means that starting time.
\newcommand{\spatialTwistCDef}[3]{
  \textcolor{spatialVelColor}{
    Ad_{\epMap{\twist{#1}{#3(t_0)} \theta_{#3}}\ldots \epMap{\twist{s}{#2-1(t_0)} \theta_{#2-1}} }\twistC{#1}{#2(t_0)}
    }
}

% #1:reference frame, #2:index of the twist, #3:index of the last link
\newcommand{\bodyTwistCDef}[3]{
  \textcolor{bodyVelColor}{
    Ad^{-1}_{\epMap{\twist{#1}{#2(t_0)} \theta_{#2}}\ldots \epMap{\twist{s}{#3(t_0)} \theta_{#3}}\transformConfig{s}{#3}{0} }\twistC{#1}{#2(t_0)}
    }
}


%%% The twist coordinate
%%% #1: notation of the joint 
\newcommand{\twistCRevolute}[1]{
\vectorTwo{
  \linearTwistASRevolute{\angularVel_{#1}}{\point{1}}
  }{
    \angularVel_{#1}
  }
}

%%% #1: notation of the joint 
\newcommand{\twistCPrismatic}[1]{
\vectorTwo{
  \linearVel_{#1}
  }{
    \zeroVector
  }
}

%%% #1:linearVel, #2: angularVel
\newcommand{\twistCDef}[2]{(#1, #2)}

%%% #1:notation of screw 
\newcommand{\twistFromScrew}[1]{
  \screwM_{#1}\vectorTwo{
    \linearTwistAS{\angularVel_{#1}}{\point{1}}{\pitch_{#1}}
  }{
    \angularVel_{#1}
  }
}

\newcommand{\twistFromScrewRevolute}[1]{
  \screwM_{#1}\vectorTwo{
    \linearTwistASRevolute{\angularVel_{#1}}{\point{1}}
  }{
    \angularVel_{#1}
  }
}

\newcommand{\vc}[1]{\bs{#1}}



\newcommand{\vectorTwo}[2]{
  \begin{bmatrix}
    #1\\
    #2
\end{bmatrix}
}

\newcommand{\vectorThree}[3]{
  \begin{bmatrix}
    #1\\
    #2\\
    #3
\end{bmatrix}
}
\newcommand{\vectorFour}[4]{
  \begin{bmatrix}
    #1\\
    #2\\
    #3\\
    #4
\end{bmatrix}
}
\newcommand{\vectorTwoRow}[2]{
  [#1^\top, #2^\top]^\top
}

\newcommand{\vectorThreeRow}[3]{
  [#1^\top, #2^\top, #3^\top]^\top
}
\newcommand{\matrixTwo}[4]{
  \begin{bmatrix}
    #1 & #2\\
    #3 & #4
\end{bmatrix}
}

\newcommand{\matrixThree}[9]{
  \begin{bmatrix}
    #1 & #2 & #3\\
    #4 & #5 & #6\\
    #7 & #8 & #9
\end{bmatrix}
}

\newcommand{\translationVel}{
  \bs{v}
}

\newcommand{\linearVel}{
\translationVel
}
\newcommand{\linearVelScalar}{
  v
}

\newcommand{\angularVel}{
  \bs{w}
}

\newcommand{\angularVelScalar}{
  w
}



%%% #1 reference frame #2 body frame
\newcommand{\spatialVel}[2]{
  \textcolor{spatialVelColor}{
    \vc{V}^{s}_{#1#2}
  }}

%%% #1 reference frame #2 body frame
\newcommand{\spatialVelDef}[2]{
  \textcolor{spatialVelColor}{
\vectorTwo{
  \spatialTVDef{#1}{#2}
  }{
    % (\rotationd{#1}{#2}\rotationInv{#1}{#2})^{\vee}
    \spatialRVDef{#1}{#2}
  }
}}

\newcommand{\spatialVelTwistDefTwo}[2]{
  \textcolor{spatialVelColor}{
    \transformdConfig{#1}{#2}{\theta}\transformInvConfig{#1}{#2}{\theta}
  }}

\newcommand{\bodyVelTwistDefTwo}[2]{
  \textcolor{spatialVelColor}{
    \transformInvConfig{#1}{#2}{\theta}\transformdConfig{#1}{#2}{\theta}
}}


% Defines the reference frame and body frame of the angular velocity
%%% #1 reference frame #2 body frame
\newcommand{\hybridVel}[2]{
  \textcolor{spatialVelColor}{
    \vc{V}^{h}_{#1#2}
  }}


\newcommand{\hybridVelDef}[2]{
  \textcolor{spatialVelColor}{
    \vectorTwo{
      \translationd{#1}{#2}
    }{
      \spatialRV{#1}{#2}
    }
  }}
\newcommand{\hybridVelDefDef}[2]{
  \textcolor{spatialVelColor}{
    \vectorTwo{
      \translationd{#1}{#2}
    }{
      \spatialRVDef{#1}{#2}
    }
  }}
%%%%%%%%%%%%%%%% Convert hybrid vel to body vel
%%% #1 Reference frame  #2 body frame
\newcommand{\hTwoB}[2]{
  \begin{bmatrix}
    \rotationInv{#1}{#2} & \zeroMatrix \\
    \zeroMatrix & \rotationInv{#1}{#2}
  \end{bmatrix}
}

%%% #1 Reference frame  #2 body frame
\newcommand{\bTwoH}[2]{
  \begin{bmatrix}
    \rotation{#1}{#2} & \zeroMatrix \\
    \zeroMatrix & \rotation{#1}{#2}
  \end{bmatrix}
}

%%%%%%%%%%%%%%%% Convert hybrid vel to body vel
\newcommand{\spatialTV}[2]{
  \textcolor{spatialVelColor}{
  \translationVel^{s}_{#1#2}
}}
\newcommand{\spatialTVDef}[2]{
  \textcolor{spatialVelColor}{
    -\rotationd{#1}{#2}\rotationInv{#1}{#2}\translation{#1}{#2} + \translationd{#1}{#2}
}}
\newcommand{\spatialTA}[2]{
  \textcolor{spatialVelColor}{
    \dot{\translationVel}^{s}_{#1#2}
  }}
\newcommand{\spatialTADef}[2]{
  \textcolor{spatialVelColor}{
    \translationdd{#1}{#2} - \spatialRATwist{#1}{#2}\translation{#1}{#2} - \spatialRVTwist{#1}{#2}\translationd{#1}{#2}
  }}
\newcommand{\spatialTADefDef}[2]{
  \textcolor{spatialVelColor}{
\translationdd{#1}{#2} -\rotationdd{#1}{#2}\rotationInv{#1}{#2} \translation{#1}{#2}  
-\rotationd{#1}{#2} \rotationInvD{#1}{#2} \translation{#1}{#2} - \rotationd{#1}{#2}\rotationInv{#1}{#2}\translationd{#1}{#2}
  }}
\newcommand{\spatialRADef}[2]{
  \textcolor{spatialVelColor}{
    \rotationdd{#1}{#2}\rotationInv{#1}{#2} + \rotationd{#1}{#2}\rotationInvD{#1}{#2}
  }}
\newcommand{\bodyRADef}[2]{
  \textcolor{bodyVelColor}{
    \rotationInv{#1}{#2}\rotationdd{#1}{#2} + \rotationInvD{#1}{#2}\rotationd{#1}{#2}
  }}
\newcommand{\spatialAcc}[2]{
  \textcolor{spatialVelColor}{
    \dot{\vc{V}}^{s}_{#1#2}
  } 
}
\newcommand{\spatialAccTwist}[2]{
  \textcolor{spatialVelColor}{
    \skewMatrix{\dot{\vc{V}}}^{s}_{#1#2}
  } 
}
\newcommand{\bodyAcc}[2]{
  \textcolor{bodyVelColor}{
    \dot{\vc{V}}^{b}_{#1#2}
  } 
}
\newcommand{\bodyAccTwist}[2]{
  \textcolor{bodyVelColor}{
    \skewMatrix{\dot{\vc{V}}}^{b}_{#1#2}
  } 
}
\newcommand{\bodyAccTwistDef}[2]{
  \textcolor{bodyVelColor}{
    \matrixTwo{
      \bodyRATwist{#1}{#2}
    }{
      \bodyTA{#1}{#2}
    }{
      \zeroVector
    }{
      0
    }
  } 
}
\newcommand{\spatialAccTwistDef}[2]{
  \textcolor{spatialVelColor}{
    \matrixTwo{
      \spatialRATwist{#1}{#2}
    }{
      \spatialTA{#1}{#2}
    }{
      \zeroVector
    }{
      0
    }
  } 
}
\newcommand{\spatialAccDef}[2]{
  \textcolor{spatialVelColor}{
    \matrixTwo{
      \spatialRATwist{#1}{#2}
    }{
      \translationdd{#1}{#2} - \spatialRATwist{#1}{#2}\translation{#1}{#2} - \spatialRVTwist{#1}{#2}\translationd{#1}{#2}
    }{
      \zeroVector
    }{
      0
    }
  }}

\newcommand{\spatialAccDefDef}[2]{
  \textcolor{spatialVelColor}{
    \matrixTwo{
      \rotationdd{#1}{#2}\rotationInv{#1}{#2} + \rotationd{#1}{#2}\rotationInvD{#1}{#2}
    }{
      \translationdd{#1}{#2} -\rotationdd{#1}{#2}\rotationInv{#1}{#2} \translation{#1}{#2}  
      -\rotationd{#1}{#2} \rotationInvD{#1}{#2} \translation{#1}{#2} - \rotationd{#1}{#2}\rotationInv{#1}{#2}\translationd{#1}{#2}
    }{
      \zeroVector
    }{
      0
    }
  }}
\newcommand{\bodyAccDefDef}[2]{
  \textcolor{bodyVelColor}{
    \matrixTwo{
      \rotationInv{#1}{#2}\rotationdd{#1}{#2} + \rotationInvD{#1}{#2}\rotationd{#1}{#2}
    }{
      \rotationInv{#1}{#2}\translationd{#1}{#2} + \rotationInv{#1}{#2}\translationdd{#1}{#2}
    }{
      \zeroVector
    }{
      0
    }
  }}
\newcommand{\bodyAccDef}[2]{
  \textcolor{bodyVelColor}{
    \matrixTwo{
      \bodyRATwist{#1}{#2}
    }{
      \bodyTV{#1}{#2} + \rotationInv{#1}{#2}\translationdd{#1}{#2}
    }{
      \zeroVector
    }{
      0
    }
  }}


\newcommand{\spatialRA}[2]{
  \textcolor{spatialVelColor}{
    \dot{\angularVel}^{s}_{#1#2}
  }}
\newcommand{\spatialRATwist}[2]{
  \textcolor{spatialVelColor}{
    \skewMatrix{\dot{\angularVel}}^{s}_{#1#2}
  }}
\newcommand{\bodyRATwist}[2]{
  \textcolor{bodyVelColor}{
    \skewMatrix{\dot{\angularVel}}^{b}_{#1#2}
  }}

\newcommand{\spatialRV}[2]{
  \textcolor{spatialVelColor}{
    \angularVel^{s}_{#1#2}
  }}

\newcommand{\spatialRVDef}[2]{
  \textcolor{spatialVelColor}{
    (\rotationd{#1}{#2}\rotationInv{#1}{#2})^{\vee}
}}
\newcommand{\spatialRVDefTwo}[2]{
  \textcolor{spatialVelColor}{
    \rotationd{#1}{#2}\rotationInv{#1}{#2}
}}

\newcommand{\spatialRVd}[2]{
  \textcolor{spatialVelColor}{
    \dot{\angularVel}^{s}_{#1#2}
}}

\newcommand{\spatialTVd}[2]{
  \textcolor{spatialVelColor}{
    \dot{\translationVel}^{s}_{#1#2}
}}

\newcommand{\spatialVelTwist}[2]{
  \textcolor{spatialVelColor}{
    \skewMatrix{\vc{V}}^{s}_{#1#2}
  }}

% \newcommand{\spatialVelTwistInv}[2]{
%   \textcolor{uncertainColor}{
%     \skewMatrix{\vc{V}}^{-s}_{#1#2}
%   }}

\newcommand{\spatialVelTwistDef}[2]{
  \textcolor{spatialVelColor}{
    \matrixTwo
    {
      \spatialRVTwist{#1}{#2}
    }{
      \spatialTV{#1}{#2}
    }{
      \zeroVector
    }{
      0
    }
}}
% \newcommand{\spatialVelTwistInvDef}[2]{
%   \textcolor{uncertainColor}{
    
% }}

\newcommand{\spatialTVTwist}[2]{
  \textcolor{spatialVelColor}{
    \skewMatrix{\vc{v}}^{s}_{#1#2}
}}
\newcommand{\spatialRVTwist}[2]{
  \textcolor{spatialVelColor}{
    \skewMatrix{\vc{w}}^{s}_{#1#2}
}}
\newcommand{\spatialRVdTwist}[2]{
  \textcolor{spatialVelColor}{
    \skewMatrix{\dot{\vc{w}}}^{s}_{#1#2}
  }}

%%% #1 Representation frame, #2 Base frame, and #3 Body frame
\newcommand{\vel}[3]{
  \textcolor{velColor}{
    \vc{V}^{#1}_{#2#3}
  }}
\newcommand{\velJacobian}[3]{
  \textcolor{velColor}{
    \jacobian^{#1}_{#2#3}
  }}
\newcommand{\tV}[3]{
  \textcolor{velColor}{
    \vc{v}^{#1}_{#2#3}
  }}

\newcommand{\rV}[3]{
  \textcolor{velColor}{
    \vc{w}^{#1}_{#2#3}
  }}


%%% #1 Reference frame #2 Body frame
\newcommand{\bodyVel}[2]{
  \textcolor{bodyVelColor}{
    \vc{V}^{b}_{#1#2}
  }}
%%% #1 Reference frame #2 Body frame
\newcommand{\bodyTVDef}[2]{
  \textcolor{bodyVelColor}{
    \rotationInv{#1}{#2}\translationd{#1}{#2}
  }}
%%% #1 Reference frame #2 Body frame
\newcommand{\bodyRVDef}[2]{
  \textcolor{bodyVelColor}{
      (\rotationInv{#1}{#2}\rotationd{#1}{#2})^{\vee}
  }}


\newcommand{\bodyRVDefTwo}[2]{
  \textcolor{bodyVelColor}{
    \rotationInv{#1}{#2}\rotationd{#1}{#2}
  }}



%%% #1 Reference frame #2 Body frame
\newcommand{\bodyVelDef}[2]{
  \textcolor{bodyVelColor}{
    \vectorTwo{
      \bodyTVDef{#1}{#2}
    }{
      \bodyRVDef{#1}{#2}
    }
  }}

%%% #1:Reference frame of the screw, #2:Body frame attached to the screw motion, #3: angle
\newcommand{\screwBodyAcc}[3]{
  \textcolor{bodyVelColor}{
    \Vee{\twistTransform{#1}{#2(0)}  \twistC{#1}{#2}}\ddot{#3}
  }}
%%% #1:Reference frame of the screw, #2:Body frame attached to the screw motion, #3: angle
\newcommand{\screwBodyVel}[3]{
  \textcolor{bodyVelColor}{
    \Vee{\twistTransform{#1}{#2(0)}  \twistC{#1}{#2}}\dot{#3}
  }}
%%% #1:Reference frame of the screw, #2:Body frame attached to the screw motion, #3: angle
\newcommand{\screwSpatialAcc}[3]{
  \textcolor{spatialVelColor}{
    \twistC{#1}{#2}\ddot{#3}
  }}

%%% #1:Reference frame of the screw, #2:Body frame attached to the screw motion, #3: angle
\newcommand{\screwSpatialVel}[3]{
  \textcolor{spatialVelColor}{
    \twistC{#1}{#2}\dot{#3}
  }}



\newcommand{\bodyVelTwist}[2]{
  \textcolor{bodyVelColor}{
    \skewMatrix{\vc{V}}^{b}_{#1#2}
  }}
% \newcommand{\bodyVelTwistInv}[2]{
%   \textcolor{uncertainColor}{
%     \skewMatrix{\vc{V}}^{-b}_{#1#2}
%   }}

\newcommand{\bodyVelTwistDef}[2]{
  \textcolor{bodyVelColor}{
    \matrixTwo
    {
      \bodyRVTwist{#1}{#2}
    }{
      \bodyTV{#1}{#2}
    }{
      \zeroVector
    }{
      0
    }
  }}
\newcommand{\bodyVelTwistDefDef}[2]{
  \textcolor{bodyVelColor}{
    \matrixTwo
    {
      \rotationInv{#1}{#2}\rotationd{#1}{#2}
    }{
      \bodyTVDef{#1}{#2}
    }{
      \zeroVector
    }{
      0
    }
}}
\newcommand{\bodyTV}[2]{
  \textcolor{bodyVelColor}{
    \vc{v}^{b}_{#1#2}
  }}
\newcommand{\bodyTA}[2]{
  \textcolor{bodyVelColor}{
    \dot{\vc{v}}^{b}_{#1#2}
  }}
\newcommand{\bodyTADef}[2]{
  \textcolor{bodyVelColor}{
    \bodyTV{#1}{#2} + \rotationInv{#1}{#2}\translationdd{#1}{#2}
  }}
\newcommand{\bodyTADefDef}[2]{
  \textcolor{bodyVelColor}{
      \rotationInv{#1}{#2}\translationd{#1}{#2} + \rotationInv{#1}{#2}\translationdd{#1}{#2}
  }}


\newcommand{\bodyRA}[2]{
  \textcolor{bodyVelColor}{
    \dot{\angularVel}^{b}_{#1#2}
  }}
% \newcommand{\bodyRADef}[2]{
%   \textcolor{bodyVelColor}{
%     \rotationInv{#1}{#2}\rotationdd{#1}{#2}
% }}
\newcommand{\bodyTVd}[2]{
  \textcolor{bodyVelColor}{
    \dot{\vc{v}}^{b}_{#1#2}
}}
\newcommand{\bodyTVTwist}[2]{
  \textcolor{bodyVelColor}{
    \skewMatrix{\vc{v}}^{b}_{#1#2}
}}

\newcommand{\bodyRV}[2]{
  \textcolor{bodyVelColor}{
    \vc{w}^{b}_{#1#2}
}}
\newcommand{\bodyRVd}[2]{
  \textcolor{geometricColorVariable}{
    \dot{\vc{w}}^{b}_{#1#2}
}}

\newcommand{\bodyForce}{
\force^{b}
}
\newcommand{\bodyTorque}{
\torque^{b}
}

\newcommand{\bodyRVTwist}[2]{
  \textcolor{bodyVelColor}{
    \skewMatrix{\vc{w}}^{b}_{#1#2}
}}



\newcommand{\jointAngle}{\theta}
\newcommand{\jointAngled}{\dot{\jointAngle}}
\newcommand{\jointAngles}{\bs{\theta}}
\newcommand{\epMap}[1]{e^{#1}}

\newcommand{\rotation}[2]{^{#2}\bs{E}_{#1}}



\newcommand{\rotationInv}[2]{R^{\top}_{#1 #2}}
\newcommand{\rotationInvD}[2]{\dot{R}^{\top}_{#1 #2}}
\newcommand{\rotationInvDD}[2]{\ddot{R}^{\top}_{#1 #2}}
\newcommand{\rotationd}[2]{\dot{R}_{#1 #2}}
\newcommand{\rotationdd}[2]{\ddot{R}_{#1 #2}}

% \newcommand{\translation}[2]{\vc{p}_{#1 #2}}
% #1 reference frame (point) #2 point
\newcommand{\translationDef}[2]{
  \textcolor{geometricColorVariable}{
    \pointer{#1}{#2} =  #2 - #1
}}


\newcommand{\translation}[2]{
  ^{#2}\vc{r}_{#1}
}

\newcommand{\translationInv}[2]{
   \textcolor{geometricColorVariable}{
     \translation{#2}{#1}
     }
}
\newcommand{\translationInvSkew}[2]{
   \textcolor{geometricColorVariable}{
     \translationSkew{#2}{#1}
     }
}
\newcommand{\translationInvDef}[2]{
   \textcolor{geometricColorVariable}{
     -\rotationInv{#1}{#2}\translation{#1}{#2}
     }
}

\newcommand{\translationSkew}[2]{\skewMatrix{\translation{#1}{#2}}}
\newcommand{\translationdSkew}[2]{\skewMatrix{\dot{\vc{p}}}_{#1 #2}}
\newcommand{\translationd}[2]{\dot{\vc{p}}_{#1 #2}}
\newcommand{\translationdd}[2]{\ddot{\vc{p}}_{#1 #2}}

% \usepackage{dsfont}
\newcommand{\identityMatrix}{\mathbf{I}} %matrix % alternativ \mathds{1} or {\mathbb{1}}
\newcommand{\idMatrix}[1]{{\identityMatrix_{#1}}} %matrix % alternativ \mathds{1} or {\mathbb{1}}

\newcommand{\zeroVector}{\mathbf{0}} %matrix % alternativ \mathds{0} or {\mathbb{0}} or \mathcal{O}
\newcommand{\zeroMatrix}{\mathbf{0}} %matrix % alternativ \mathds{0} or {\mathbb{0}} or \mathcal{O}

\newcommand{\SE}[1]{SE{(#1)}}
\newcommand{\se}[1]{se{(#1)}}
\newcommand{\ssRegion}{{\set{\text{ssr}}}}
\newcommand{\ssrUB}{{l_1}}
\newcommand{\ssrLB}{{l_2}}
\newcommand{\SO}[1]{SO{(#1)}}
%\newcommand{\so}[1]{so{(#1)}}
\newcommand{\soPowerTwoDef}[1]{
  \projectionDef{#1} - \twoNorm{#1} \identityMatrix
}
\newcommand{\soPowerThreeDef}[1]{
  - \twoNorm{#1} \skewMatrix{#1}
}


\newcommand{\transform}[2]{
  \textcolor{geometricColorVariable}{
    g_{#1#2}
  }}

\newcommand{\transformd}[2]{
 \textcolor{geometricColorVariable}{
   \dot{g}_{#1#2}
 }}
\newcommand{\transformdDef}[2]{
 \textcolor{geometricColorVariable}{
   \matrixTwo{
     \rotationd{#1}{#2}
   }{
     \translationd{#1}{#2}
   }{
     \zeroVector
   }{
     0
   }
 }}
\newcommand{\transformdd}[2]{
 \textcolor{geometricColorVariable}{
   \ddot{g}_{#1#2}
 }}
\newcommand{\transformddDef}[2]{
 \textcolor{geometricColorVariable}{
   \matrixTwo{
     \rotationdd{#1}{#2}
   }{
     \translationdd{#1}{#2}
   }{
     \zeroVector
   }{
     0
   }
 }}
\newcommand{\transformInv}[2]{
  \textcolor{geometricColorVariable}{
    g^{-1}_{#1#2}
  }}
\newcommand{\transformInvD}[2]{
  \textcolor{geometricColorVariable}{
    % \dot{g}^{-1}_{#1#2}
    \derivative{g^{-1}_{#1#2}}
  }}
\newcommand{\transformInvDConfig}[3]{
  \textcolor{geometricColorVariable}{
    % \dot{g}^{-1}_{#1#2}
    \derivative{g^{-1}_{#1#2}(#3)}
  }}
\newcommand{\transformInvDDef}[2]{
  \textcolor{geometricColorVariable}{
    \matrixTwo{
      \rotationInvD{#1}{#2}
    }{
      -\rotationInvD{#1}{#2} \translation{#1}{#2}  - \rotationInv{#1}{#2} \translationd{#1}{#2}
    }{
      \zeroVector
    }{
      0
    }
  }}

\newcommand{\transformDef}[2]{
  \textcolor{geometricColorVariable}{
    \matrixTwo{
      \rotation{#1}{#2}
    }{
      \translation{#1}{#2}
    }{
      \zeroVector
    }{
      1
    }
  }
}

\newcommand{\transformInvDef}[2]{
  \textcolor{geometricColorVariable}{
    \matrixTwo{
      \rotationInv{#1}{#2}
    }{
      -\rotationInv{#1}{#2} \translation{#1}{#2} 
    }{
      \zeroVector
    }{
      1
    }
  }
}
\newcommand{\mapDef}[2]{
  : #1 \rightarrow #2
}


% transform a twist coordinate from coordinate #1 to coordinate #2
\newcommand{\twistCTransform}[2]{
  \textcolor{geometricColorVariable}{
    \twistTransform{#1}{#2}
  }
}
\newcommand{\twistCTransformDef}[2]{
  \textcolor{geometricColorVariable}{
    \twistTransformDef{#1}{#2}
  }
}
% transform a point from coordinate #1 to coordinate #2
\newcommand{\pointTransform}[2]{
  \textcolor{geometricColorVariable}{
    \transformInv{#1}{#2}
  }
}
\newcommand{\pointTransformDef}[2]{
  \textcolor{geometricColorVariable}{
    \transformInvDef{#1}{#2}
  }
}
\newcommand{\pointTransformTwo}[2]{
  \textcolor{geometricColorVariable}{
    \transform{#2}{#1}
  }
}
\newcommand{\pointTransformTwoDef}[2]{
  \textcolor{geometricColorVariable}{
    \transformDef{#2}{#1}
  }
}
% transform from coordinate #1 to coordinate #2
\newcommand{\twistTransformTwo}[2]{
  \textcolor{geometricColorVariable}{
    \adg{#2}{#1}
  }
}

\newcommand{\twistTransformdTwo}[2]{
  \textcolor{geometricColorVariable}{
    \derivative{\adg{#2}{#1}}
  }
}
\newcommand{\twistTransformdTwoDef}[2]{
  \textcolor{geometricColorVariable}{
    -\adg{#2}{#1}(\spatialVel{#2}{#1}\times)
  }
}

\newcommand{\twistTransformTwoDef}[2]{
  \textcolor{geometricColorVariable}{
    \adgDef{#2}{#1}
  }
}

\newcommand{\twistTransform}[2]{
  \textcolor{geometricColorVariable}{
    \adgInv{#1}{#2}
    }
  }
  

\newcommand{\twistTransformDef}[2]{
  \adgInvDef{#1}{#2}
}
\newcommand{\twistTransformd}[2]{
  \textcolor{geometricColorVariable}{
    \derivative{\twistTransform{#1}{#2}}
    }
}
\newcommand{\twistTransformdDef}[2]{
  \textcolor{geometricColorVariable}{
    -\twistTransform{#1}{#2}(\spatialVel{#1}{#2}\times)
    }
}




% #1: reference frame , #2 link(body) frame, and #3: New reference frame.
\newcommand{\spatialVelTransform}[3]{
  \textcolor{spatialVelColor}{
    % \spatialVel{#3}{#1} + \spatialVelStaticTransform{#3}{#1}\spatialVel{#1}{#2}
    \spatialVel{#3}{#1} + \twistTransform{#1}{#3}\spatialVel{#1}{#2}
  }
}
% #1: reference frame , #2 link(body) frame, and #3: New reference frame.
\newcommand{\spatialAccTransform}[3]{
  \textcolor{spatialVelColor}{
    % \spatialVel{#3}{#1} + \spatialVelStaticTransform{#3}{#1}\spatialVel{#1}{#2}
    \spatialAcc{#3}{#1} 
    +\twistTransform{#1}{#3} (\spatialVelTwist{#1}{#2}\spatialVelTwist{#1}{#3} - \spatialVelTwist{#1}{#3}\spatialVelTwist{#1}{#2})^\vee
    % +\twistTransform{#1}{#3} (\bodyVelTwist{#3}{#1}\bodyVelTwist{#2}{#1})^\vee
    + \twistTransform{#1}{#3}\spatialAcc{#1}{#2} 
  }
}
% #1: reference frame , #2 link(body) frame, and #3: New reference frame.
\newcommand{\spatialAccTransformTwo}[3]{
  \textcolor{spatialVelColor}{
    \spatialAcc{#3}{#1}
    +\twistTransform{#1}{#3}\cross{\spatialVel{#1}{#2}}{\spatialVel{#1}{#3}}
    + \twistTransform{#1}{#3}\spatialAcc{#1}{#2} 
  }
}

% #1: reference frame , #2 link(body) frame, and #3: New reference frame.
% \newcommand{\bodyAccTransform}[3]{
%   \textcolor{spatialVelColor}{
%     \bodyAcc{#1}{#2}
%     + 2\twistTransform{#1}{#2}(
%     \bodyVelTwist{#3}{#1}\spatialVelTwist{#1}{#2}
%     )^{\vee}
%     + \twistTransform{#1}{#2}(
%     \spatialVelTwist{#1}{#3}\spatialVelTwist{#1}{#2}
%     )^{\vee}
%     + \twistTransform{#1}{#2}(
%     \bodyVelTwist{#2}{#1}\bodyVelTwist{#3}{#1}
%     )^{\vee}
%     +\twistTransform{#1}{#2}\bodyAcc{#3}{#1}
%   }
%   }
% #1: reference frame , #2 link(body) frame, and #3: New reference frame.
\newcommand{\bodyAccTransform}[3]{
  \textcolor{bodyVelColor}{
    \bodyAcc{#1}{#2}
    + \twistTransform{#1}{#2}(
    \bodyVelTwist{#2}{#1}\bodyVelTwist{#3}{#1}
    -\bodyVelTwist{#3}{#1}\bodyVelTwist{#2}{#1}
    )^{\vee}
    +\twistTransform{#1}{#2}\bodyAcc{#3}{#1}
  }
}
% #1: reference frame , #2 link(body) frame, and #3: New reference frame.
\newcommand{\bodyAccTransformTwo}[3]{
  \textcolor{bodyVelColor}{
    \bodyAcc{#1}{#2}
    % \twistTransformdDef{#1}{#2}
    +\twistTransform{#1}{#2}\bodyVel{#2}{#1}\times
    \bodyVel{#3}{#1}
    +\twistTransform{#1}{#2}\bodyAcc{#3}{#1}
  }
}

\newcommand{\spatialVelTransformTwo}[3]{
  \textcolor{spatialVelColor}{
    % \spatialVel{#3}{#1} + \spatialVelStaticTransform{#3}{#1}\spatialVel{#1}{#2}
    \spatialVel{#3}{#1} +
    \twistTransformTwo{#1}{#3}\spatialVel{#1}{#2}
  }
}

% Expressed in #2
% #1: reference frame , #2 link(body) frame, and #3: New reference frame.
\newcommand{\bodyVelTransform}[3]{
  \textcolor{bodyVelColor}{
    \twistTransform{#1}{#2}\bodyVel{#3}{#1} + \bodyVel{#1}{#2}
  }
}
\newcommand{\bodyVelTransformTwo}[3]{
  \textcolor{bodyVelColor}{
    \twistTransformTwo{#1}{#2}\bodyVel{#3}{#1} + \bodyVel{#1}{#2}
  }
}


% % transform body velocity in frame #1 to spatial velocity expressed in frame #2:
% % #1 body frame, #2 spatial frame
% \newcommand{\bodyVelToSpatialVel}[2]{
% \adg{#2}{#1}
% }
% \newcommand{\bodyVelToSpatialVelDef}[2]{
% \adgDef{#2}{#1}
% }

% % transform spatial velocity expressed in frame # 1 to body frame #2
% % #1 spatial frame, #2 body frame
% \newcommand{\spatialVelToBodyVel}[2]{
% \adgInv{#1}{#2}
% }
% \newcommand{\spatialVelToBodyVelDef}[2]{
% \adgInvDef{#1}{#2}
% }


% Transform from #1 to #2 
\newcommand{\svaTransform}[2]{
  \textcolor{svaColorVariable}{
    {^{#2}T_{#1}}
  }
}
% Transform from #1 to #2 
\newcommand{\svaTransformDef}[2]{
  \textcolor{svaColorVariable}{
  \matrixTwo{
    \rotation{#2}{#1}
  }{
    -\rotation{#2}{#1}\translation{#1}{#2}
  }{
    \zeroVector
  }{
    1
  }
  }
}
\newcommand{\svaTransformDefTwo}[2]{
  \textcolor{svaColorVariable}{
  \matrixTwo{
    \rotationInv{#1}{#2}
  }{
    -\rotationInv{#1}{#2}\translation{#1}{#2}
  }{
    \zeroVector
  }{
    1
  }
  }
}

\newcommand{\svaT}{
  \textcolor{svaColorVariable}{
    \bs{X}
  }
}

% Transform from #1 to #2 (#2 is represented in #1)
\newcommand{\svaVT}[2]{
  \textcolor{svaColorVariable}{
  {^{#2}\svaT_{#1}}
  % \svaT{#1}{#2}
  }
}

% Transform Vel from #1 to #2 
\newcommand{\svaVTDef}[2]{
  % eq.~2.24 of FeatherStone
  % \matrixTwo{
  %   \rotationInv{#1}{#2}
  % }{\zeroMatrix}{
  %   -\rotationInv{#1}{#2}\translationSkew{#1}{#2}
  % }
  % {
  %   \rotationInv{#1}{#2}
  % }
  \textcolor{svaColorVariable}{
  \matrixTwo{
    \rotation{#2}{#1}
  }{\zeroMatrix}{
    -\rotation{#2}{#1}\translationSkew{#1}{#2}
  }
  {
    \rotation{#2}{#1}
  }
  }
}

% Transform Vel from #1 to #2 
\newcommand{\svaVTDefTwo}[2]{
  % eq.~2.26 of FeatherStone
  \textcolor{svaColorVariable}{
  \matrixTwo{
    \rotationInv{#1}{#2}
  }{\zeroMatrix
  }{
    \translationSkew{#2}{#1}\rotationInv{#1}{#2}
  }
  {
    \rotationInv{#1}{#2}
  }
  }
}

% Transform from #1 to #2 
\newcommand{\svaFT}[2]{
  \textcolor{svaColorVariable}{
    {^{#2}\svaT^{*}_{#1}}
  }
}

% Transform from #1 to #2 
% \newcommand{\svaFTDef}[2]{
\newcommand{\svaFTDef}[2]{
  % Following eq.2.27 of Featherstone
  % \matrixTwo
  % {\rotationInv{#2}{#1}}
  % {
  %   \translationSkew{#1}{#2}\rotationInv{#2}{#1}
  % }{
  %   \zeroMatrix
  % }
  %   {\rotationInv{#2}{#1}}
  \textcolor{svaColorVariable}{
  \matrixTwo
  {\rotationInv{#1}{#2}}
  {
    \translationSkew{#2}{#1}\rotationInv{#1}{#2}
  }{
    \zeroMatrix
  }
  {\rotationInv{#1}{#2}}
  }
}

% Transform from #1 to #2 
\newcommand{\svaFTDefTwo}[2]{
  % Following eq.2.25 of Featherstone
  % \matrixTwo{
  %   \rotationInv{#1}{#2}
  % }{
  %   -\rotationInv{#1}{#2}\translationSkew{#1}{#2}
  % }{
  %   \zeroMatrix
  % }
  %   {\rotationInv{#1}{#2}}
  \textcolor{svaColorVariable}{
  \matrixTwo{
    \rotation{#1}{#2}
  }{
    -\rotation{#1}{#2}\translationSkew{#1}{#2}
  }{
    \zeroMatrix
  }
  {\rotation{#1}{#2}}
  }
}

% Reference frame: #1, link frame: #2 
\newcommand{\adg}[2]{
  \textcolor{geometricColorVariable}{
    Ad_{g_{#1#2}}
  }
}

\newcommand{\adgDef}[2]{
  \textcolor{geometricColorVariable}{
  \begin{bmatrix}
    \rotation{#1}{#2} & \translationSkew{#1}{#2}\rotation{#1}{#2} \\
    \zeroMatrix & \rotation{#1}{#2}
  \end{bmatrix}
}}

\newcommand{\adgInv}[2]{
  \textcolor{geometricColorVariable}{
  Ad^{-1}_{g_{#1#2}}
}}
\newcommand{\adgTwoInv}[2]{
  \textcolor{geometricColorVariable}{
  Ad_{g^{-1}_{#1#2}}
}}

\newcommand{\adgTrans}[2]{
  \textcolor{geometricColorVariable}{
  Ad^{\top}_{g_{#1#2}}
}}
\newcommand{\adgTransDef}[2]{
  \textcolor{geometricColorVariable}{
  \begin{bmatrix}
    \rotationInv{#1}{#2} & \zeroMatrix \\
     -\rotationInv{#1}{#2}\translationSkew{#1}{#2} & \rotationInv{#1}{#2}
\end{bmatrix}
}}


\newcommand{\adgInvDefTwo}[2]{
  \textcolor{geometricColorVariable}{
  \begin{bmatrix}
    \rotationInv{#1}{#2} & -(\rotationInv{#1}{#2}\translation{#1}{#2})^{\skewMatrix{}}\rotationInv{#1}{#2} \\
    \zeroMatrix & \rotationInv{#1}{#2}
\end{bmatrix}
}}
\newcommand{\adgInvDef}[2]{
  \textcolor{geometricColorVariable}{
  \begin{bmatrix}
    \rotationInv{#1}{#2} & -\rotationInv{#1}{#2}\translationSkew{#1}{#2} \\
    \zeroMatrix & \rotationInv{#1}{#2}
\end{bmatrix}
}}

% From frame #1 to frame #2 
\newcommand{\geometricFTDef}[2]{
  \textcolor{geometricColorVariable}{
    \adgInvTransDef{#2}{#1}
  }
%   \begin{bmatrix}
%     \rotation{#1}{#2} & \zeroMatrix \\
%     \translationSkew{#1}{#2}\rotation{#1}{#2}  & \rotation{#1}{#2}
% \end{bmatrix}
}

% From frame #1 to frame #2 
\newcommand{\geometricFTTwo}[2]{
  \textcolor{geometricColorVariable}{
    \adgTrans{#1}{#2}
  }
}
\newcommand{\geometricFTTwoDef}[2]{
  \textcolor{geometricColorVariable}{
    \adgTransDef{#1}{#2}
  }
}

% From frame #1 to frame #2 
\newcommand{\geometricFT}[2]{
  \textcolor{geometricColorVariable}{
    \adgInvTrans{#2}{#1}
  }}

\newcommand{\adgInvTrans}[2]{
  \textcolor{geometricColorVariable}{
  Ad^{\top}_{g^{-1}_{#1#2}}
}}

\newcommand{\adgInvTransDef}[2]{
  \textcolor{geometricColorVariable}{
  \begin{bmatrix}
    \rotation{#1}{#2} & \zeroMatrix \\
    \translationSkew{#1}{#2}\rotation{#1}{#2}  & \rotation{#1}{#2}
  \end{bmatrix}
  }
}



\newcommand{\adgTransform}[3]{
  \textcolor{geometricColorVariable}{
  \transform{#1}{#2} \skewMatrix{#3} \transformInv{#1}{#2}
}}
\newcommand{\wrench}{\bs{w}}
% The contact wrench cone:
\newcommand{\cwc}{\text{C}}
\newcommand{\cwcFrame}[2]{{^{#1}\cwc_{#2}}}

\newcommand{\bodyWrench}{\bs{W}^{b}}
\newcommand{\wrenchFrame}[2]{{^{#1}\wrench_{#2}}} % The arguments are: (1)Frame (2) Application point
\newcommand{\cone}[1]{{\text{K}_{#1}}}
\newcommand{\generator}[1]{{\bs{\lambda}_{#1}}}
\newcommand{\generatorScalar}[1]{{\lambda_{#1}}}

\newcommand{\frictionCone}{{^{\coefF}K}}
\newcommand{\normalCone}{{^{\surfaceNormal}K}}
\newcommand{\impulseGenerator}{\bs{\lambda}}


%%%%%%%%%%%%%%%%%%%%%%%%% Inertia and momentum:
% #1 body frame name, the COM and body inertia are calculated in its own coordinate frame
% \newcommand{\bodyGInertia}[1]{
%   \body{\gInertia}_{#1}
% }
%   #1 reference frame name, #2 body frame name
%  locally, we can specify: #1 body frame: F_i  #2 the local centroidal frame: F_com_i
\newcommand{\bodyGInertia}[2]{
\textcolor{bodyVelColor}{
  \body{\gInertia}_{#1 #2}
  }
}


\newcommand{\mechanicalConnection}{
  \bs{A}
}




% #1 reference frame name, #2 body frame name, the COM and body inertia are calculated in its own coordinate frame
% \newcommand{\bodyGInertiaDef}[1]{
%   \matrixTwo{\mass_{#1} \identityMatrix }{\zeroMatrix}{\zeroMatrix}{\bodyMetaInertia{#1}}
% }

\newcommand{\bodyGInertiaDef}[2]{
  \textcolor{bodyVelColor}{
    \matrixTwo{\mass_{#2} \identityMatrix }{\zeroMatrix}{\zeroMatrix}{\bodyMetaInertia{#1}{#2}}
  }
}

% \newcommand{\bodyMetaInertiaDef}[2]{
%   \rotation 
% }
\newcommand{\svaBodyGInertiaDef}[2]{
  \textcolor{svaColorVariable}{
    \matrixTwo{\bodyMetaInertia{#1}{#2} }{\zeroMatrix}{\zeroMatrix}{\mass_{#2} \identityMatrix}
  }
}


% #1 Reference frame #2 body frame
\newcommand{\spatialGInertia}[2]{
  \textcolor{spatialVelColor}{
    \spatial{\gInertia}_{#1#2}
  }
}

\newcommand{\spatialGInertiad}[2]{
  \textcolor{spatialVelColor}{
    \spatial{\dot{\gInertia}}_{#1#2}
  }
}

\newcommand{\spatialGInertiadDef}[2]{
  \textcolor{spatialVelColor}{
    -\transpose{(\cross{\spatialVel{#1}{#2}}{})} \spatialGInertia{#1}{#2} - \spatialGInertia{#1}{#2}\cross{\spatialVel{#1}{#2}}{}
  }
}

% #1 Current reference frame #2 New Reference frame
\newcommand{\momentaTransform}[2]{
  \transpose{\twistTransform{#2}{#1}}
}
\newcommand{\momentaTransformDef}[2]{
  \transpose{\twistTransformDef{#2}{#1}}
}

% #1 Current reference frame #2 New Reference frame
\newcommand{\svaMomentaTransform}[2]{
  \transpose{\svaVT{#2}{#1}}
}
\newcommand{\svaMomentaTransformDef}[2]{
  \transpose{\svaVTDef{#2}{#1}}
}


% #1 and #2 defines the spatial generalized inertia to be transformed.
% #1 Reference frame #2 body frame 
% #3 defines the new reference frame
\newcommand{\spatialGInertiaTransform}[3]{
  \transpose{\twistTransform{#3}{#1}}\spatialGInertia{#1}{#2}\twistTransform{#3}{#1}
}



\newcommand{\svaSpatialGInertiaTransform}[3]{
  \transpose{\svaVT{#3}{#1}}\spatialGInertia{#1}{#2}\svaVT{#3}{#1}
}

%%% Generalized inertia transform
% #1:new reference frame
% #2:current reference frame
% #3:current Generalized inertia
\newcommand{\gInertiaTransform}[3]{
  \transpose{\twistTransform{#1}{#2}}
  #3
  \twistTransform{#1}{#2}
}

\newcommand{\svaGInertiaTransform}[3]{
  \transpose{\svaVT{#1}{#2}}
  #3
  \svaVT{#1}{#2}
}


% #1 Reference frame #2 body frame
\newcommand{\spatialGInertiaDef}[2]{
  \transpose{\twistTransform{#1}{#2}}
  % \bodyGInertia{#1}{#2}
  \spatialGInertia{#2}{\com_{#2}}
  \twistTransform{#1}{#2}
}

\newcommand{\spatialGInertiaDefTwo}[2]{
  \transpose{\twistTransformTwo{#1}{#2}}
  % \bodyGInertia{#1}{#2}
  \spatialGInertia{#2}{\com_{#2}}
  \twistTransformTwo{#1}{#2}
}
% #1 Reference frame #2 body frame
\newcommand{\spatialGInertiaDefDef}[2]{
  % \transpose{{\adgInvDef{#1}{#2}}} \bodyGInertiaDef{#2}\adgInvDef{#1}{#2}
  \transpose{{\twistTransformDef{#1}{#2}}}
  \spatialGInertia{#2}{\com_{#2}}
  % \bodyGInertiaDef{#1}{#2}
  \twistTransformDef{#1}{#2}
}
\newcommand{\spatialGInertiaDefDefTwo}[2]{
  % \transpose{{\adgInvDef{#1}{#2}}} \bodyGInertiaDef{#2}\adgInvDef{#1}{#2}
  \transpose{{\twistTransformTwoDef{#1}{#2}}}
  % \bodyGInertiaDef{#1}{#2}
  \spatialGInertia{#2}{\com_{#2}}
  \twistTransformTwoDef{#1}{#2}
}

% #1 Reference frame #2 body frame
\newcommand{\spatialGInertiaDetail}[2]{
  \textcolor{spatialVelColor}{
\matrixTwo{
  \mass_{#2} \identityMatrix
}{
  \mass_{#2}\transpose{\translationSkew{#1}{#2}}
}{
  \mass\translationSkew{#1}{#2}
}{
  \mass \translationSkew{#1}{#2} \transpose{\translationSkew{#1}{#2}} + \rotation{#1}{#2}\bodyMetaInertia{#1}{#2}\rotationInv{#1}{#2}
}
}
}
% #1 Reference frame #2 body frame
\newcommand{\spatialGInertiaDiagonalDetail}[2]{
  \textcolor{spatialVelColor}{
\matrixTwo{
  \mass_{#2} \identityMatrix
}{
  \zeroMatrix
}{
  \zeroMatrix
}{
  \rotation{#1}{#2}\bodyMetaInertia{#1}{#2}\rotationInv{#1}{#2}
}
}
}

% #1 Reference frame #2 body frame
\newcommand{\svaSpatialGInertia}[2]{
  \textcolor{svaColorVariable}{
    \transpose{\svaVT{#1}{#2}} \bodyGInertia{#1}{#2}\svaVT{#1}{#2}
  }
}

% #1 Reference frame #2 body frame
\newcommand{\svaSpatialGInertiaDef}[2]{
  \textcolor{svaColorVariable}{
    \transpose{\svaVTDef{#1}{#2}} \svaBodyGInertiaDef{#1}{#2}\svaVTDef{#1}{#2}
  }
}


\newcommand{\metaInertia}{\mathcal{I}}



% #1: the reference frame and #2 the link frame
\newcommand{\spatialMetaInertia}[2]
{
\textcolor{spatialVelColor}{ 
  \spatial{\metaInertia}_{{#1}{#2}}
}
}
% #1: the reference frame and #2 the link frame
\newcommand{\bodyMetaInertia}[2]
{
  \textcolor{bodyVelColor}{ 
    \body{\metaInertia}_{{#1 #2}}
    }
}
\newcommand{\spatialInertia}{\mathcal{I}^{\prime}}

\newcommand{\gInertia}{\mathcal{M}}

\newcommand{\kineticEnergy}{T(\jointPosition, \jointPositiond)}
\newcommand{\potentialEnergy}{V(\jointPosition)}
\newcommand{\totalEnergy}{H}
\newcommand{\lagR}{\mathcal{L}(\jointPosition, \jointPositiond)}

\newcommand{\lagRDef}{
  \quadratic{\jointPositiond}{\inertiaMatrix} + \potentialEnergy
}

%%% #1: dof number
\newcommand{\lagRDefDef}[1]{
  \frac{1}{2}\sum^{#1}_{i,j=1}\inertiaMatrix_{ij}(\jointPosition)\jointAngled_i\jointAngled_j - \potentialEnergy
}

% The Lagrange's equation of a link
%%% #1: link index
\newcommand{\lagREquations}[1]{
  \derivative{\pd{\lagR}{\jointAngled_{#1}}} - \pd{\lagR}{\jointPosition_{#1}} = \torque_{#1}
}

% The Lagrange's equation of a link
%%% #1: link index
\newcommand{\lagREquationsDef}[1]{
  \agg{j}{\dof}{
  (\inertiaMatrixd_{#1j}\jointPositiond_j  + \inertiaMatrix_{#1j}\jointPositiondd_j)
}
- \frac{1}{2}\agg{#1}{\dof}
{
  \pd{\inertiaMatrix_{#1j}}{\jointPosition_#1} \jointPositiond_#1\jointPositiond_j
  - \pd{\potentialEnergy}{\jointPosition_#1}
}
  = \torque_{#1}
}


\newcommand{\lagREquationsDefDef}[1]{
  \agg{j}{\dof}{
    \inertiaMatrix_{#1j}\jointPositiondd_j
  }
  + \frac{1}{2}\agg{j,k}{\dof}{(
    \pd{\inertiaMatrix_{#1j}}{\jointPosition_k}
    +\pd{\inertiaMatrix_{#1k}}{\jointPosition_j}
    - \pd{\inertiaMatrix_{kj}}{\jointPosition_#1}
    % \christoffelSymbolDef{#1}{j}{k}
  )\jointPositiond_k\jointPositiond_j}
+ \pd{\potentialEnergy}{\jointPosition_#1}
  = \torque_{#1}
}
%%%% #1 link index #2,#3 iteration index
\newcommand{\christoffelSymbol}[3]{
  \Gamma_{#1#2#3}
}
\newcommand{\christoffelSymbolDef}[3]{
  \frac{1}{2}(
    \pd{\inertiaMatrix_{#1#2}}{\jointPosition_#3}
    +\pd{\inertiaMatrix_{#1#3}}{\jointPosition_#2}
    - \pd{\inertiaMatrix_{#3#2}}{\jointPosition_#1})
}


% The argument are: force, torque
\newcommand{\wrenchC}[2]{
\begin{bmatrix}
#1\\
#2
\end{bmatrix}
}

% The argument are: force, torque
\newcommand{\wrenchR}[2]{
[#1^\top, #2^\top]^\top
}

\newcommand{\bodyJacobian}[2]{
  \textcolor{bodyVelColor}{
    \jacobian^{b}_{#1 #2}
  }
}

\newcommand{\bodyJacobiand}[2]{
  \textcolor{bodyVelColor}{
    \dot{\jacobian}^{b}_{#1 #2}
  }
}


%%%%%%%%%%%%%%%%%%%%%%%%%% QP Control law

% #1 velocity #2 mass
\newcommand{\quadraticObj}[2]{\transpose{#1}#2#1}

\newcommand{\desired}[1]{#1^{*}}
\newcommand{\reff}[1]{#1^{\text{ref}}}

% integration error 
\newcommand{\iError}[1]{\int{\bs{e}_{#1}}}
\newcommand{\iErrorDef}[1]{\int{\errorDef{#1}}dt}

\newcommand{\error}[1]{\bs{e}_{#1}}
\newcommand{\errord}[1]{\dot{\bs{e}}_{#1}}
\newcommand{\errordd}[1]{\ddot{\bs{e}}_{#1}}

\newcommand{\errorDef}[1]{
  #1 - \reff{#1}
}
% #1:index of the input
\newcommand{\controlInput}[1]{u_{#1}}

%% #1 real part, #2 imaginery part
\newcommand{\complexNum}[2]{
  (#1 ~ #2*i)
}

%% #1 upper bound, #2 middle #3 lower bound
\newcommand{\saturation}[3]{
\left\{
  \begin{array}{lr}
    #1 \\
    #2 \\
    #3
  \end{array}
\right.
}
% #1 velocity #2 mass
\newcommand{\quadratic}[2]{\frac{1}{2} #1^\top #2 #1}
\newcommand{\objF}{J}


%%%%%%%%%%%%%%%%%%%%%%%%%% Cross product 
\newcommand{\paral}[2]{#1 \parallel #2}
\newcommand{\perpen}[2]{#1 \perp #2}

\newcommand{\bs}{\boldsymbol}
\newcommand{\backfill}{\hfill\(\blacksquare\)}

\newcommand{\skewMatrix}[1]{#1\times}

\newcommand{\skewMatrixTwo}[1]{({#1})^{\widehat{~}}}


\newcommand{\vectriple}[3]{#1\times(#2 \times #3)} 
\newcommand{\vectripleexpan}[3]{
(\inner{#1}{#3}){#2} - (\innerP{#1}{#2}){#3}
}
\newcommand{\vectripleexpannegate}[3]{
-(\inner{#1}{#3}){#2} + (\innerP{#1}{#2}){#3}
} 
\newcommand{\scalartriple}[3]{{#1}\bullet({#2} \times {#3})}
\newcommand{\vcd}[3]{({#1}\bullet{#3}){#2} - ({#1}\bullet{#2}) {#3}}  % Vector cross product definition  

\newcommand{\inner}[2]{#1\cdot#2} % Inner product 
\newcommand{\innerDef}[2]{\transpose{{#1}}#2} % Inner product 
\newcommand{\normal}[1]{\overline{{#1}}}

\newcommand{\nc}[1]{\frac{#1}{| #1| }}% normal calc

\newcommand{\transpose}[1]{{#1}^\top}
\newcommand{\transInverse}[1]{#1^{-\top}}
%% \newcommand{\pseudoInverse}[1]{#1^{\dagger}}
%% \newcommand{\inverse}[1]{{#1}^{-1}}


\newcommand{\pseudoInverse}[1]{{#1}^{+}}
% When  #row >  #column
\newcommand{\pseudoInverseColumnDef}[1]{\inverse{( \transpose{#1} #1 )}\transpose{#1}}
% When  #column >  #row
\newcommand{\pseudoInverseRowDef}[1]{\transpose{#1}\inverse{( #1 \transpose{#1} )}}

\newcommand{\weightedPseudoInverse}[1]{#1^{\dagger}}
% % #1 Jacobian #2 The weight, e.g., the inertia matrix. 
\newcommand{\weightedPseudoInverseRowDef}[2]{\inverse{#2}\transpose{#1}\inverse{( #1 \inverse{#2} \transpose{#1} )}}
\newcommand{\weightedPseudoInverseColumnDef}[2]{\inverse{( \transpose{#1} \inverse{#2} #1 )}\transpose{#1}\inverse{#2}}
\newcommand{\inverse}[1]{{#1}^{-1}}

%%%%%%%%%%%%%%%%%%%%%%%%%% Resultant torque
% The argument are: force, torque, application point of the input wrench, application point of the resultant wrench. 
\newcommand{\resultantTorque}[4]{#2 + (#3 - #4)\times #1}
\newcommand{\resultantTorqueTwo}[4]{#2 + \pointer{#4}{#3}\times #1}
\newcommand{\torqueFrame}[2]{{\torque_{#1#2}}}



% \boldsymbol{#1} + (\boldsymbol{#2} - \boldsymbol{#3})\times \boldsymbol{#4}}

% The argument is frame name. 
\newcommand{\cframe}[1]{\mathcal{F}_{#1}}

% Denotes the origin of a coordinate frame
\newcommand{\origin}[1]{#1_{o}}
\newcommand{\todo}[1]{To do:~
\textcolor{bodyVelColor}{
  #1
  }
}
% \newcommand{\wrench}{\bs{W}} %TODO choose a small character because it is a vector


\newcommand{\numEE}{m}

\newcommand{\forceJacobian}[2]{\jacobian_{#1 #2}}

\newcommand{\configurationFrame}[2]{g_{#1 #2}}



\newcommand{\bodyFrame}[1]{{^b#1}}

\newcommand{\spatialJacobian}[2]{
  \textcolor{spatialVelColor}{
    \spatial{\jacobian}_{#1 #2}
  }
}

\newcommand{\spatialJacobiand}[2]{
  \textcolor{spatialVelColor}{
    \dot{\spatial{\jacobian}}_{#1 #2}
  }
}

\newcommand{\forceJacobianDef}[2]{
\begin{bmatrix}
  \rotation{#1}{#2} \\
  \pointer{#1}{#2} \times \rotation{#1}{#2}
\end{bmatrix}}

% Transform from #1 to #2 
\newcommand{\svaForceJacobianDef}[2]{
  \textcolor{svaColorVariable}{  
    \vectorTwo{
      \translationSkew{#2}{#1}\rotationInv{#1}{#2}    
    }{
      {\rotationInv{#1}{#2}}
    }
  }
}
% Transform from #1 to #2 
\newcommand{\svaForceJacobianDefTwo}[2]{
  \textcolor{svaColorVariable}{
    \vectorTwo{
      -\rotation{#2}{#1}\translationSkew{#1}{#2}
    }{
      {\rotation{#2}{#1}}
    }
  }
}


\newcommand{\mass}{\text{m}}

\newcommand{\numContacts}{{\text{m}_{\text{sc}}}}
\newcommand{\numImpacts}{{m_2}}
\newcommand{\numFree}{\numEE - \numContacts -\numImpacts }
\newcommand{\contactPoint}{\bs{p}}

\newcommand{\ee}{\text{e}}



% The argument is point name. 
\newcommand{\point}[1]{\boldsymbol{q}_{#1}}

\newcommand{\force}{\bs{f}}
\newcommand{\forceScalar}{f}
% The arguments are: (1) frame (2) Application point. 
% \newcommand{\forceFrame}[2]{{^{#1}\bs{f}_{#2}}}
\newcommand{\forceFrame}[2]{{^{#1}\force_{#2}}}
\newcommand{\impulseFrame}[2]{{\impulse_{#1#2}}}

% The arguments are: (1) Equivalent frame name (2) Application point of the force 
\newcommand{\forceGraspMatrix}[2]{
\begin{bmatrix}
    \rotation{\bs{#1}}{#2}^\top \\
    \pointer{#1}{#2} \times \rotation{#1}{#2}^\top 
  \end{bmatrix}
}
\newcommand{\graspMatrix}[2]{
  {G_{#1#2}}
}

\newcommand{\graspMatrixDef}[2]{
\begin{bmatrix}
  \rotation{#1}{#2} & 0\\
  \pointer{#1}{#2} \times \rotation{#1}{#2} & \rotation{#1}{#2}
\end{bmatrix}
}
\newcommand{\comPlane}{\com_{x,y}}
\newcommand{\pendulumC}{\text{w}}
\newcommand{\pendulumCDef}{\sqrt{\frac{g}{\com_z}}}
\newcommand{\comVel}{\dot{\com}}
\newcommand{\comVelPlane}{\dot{\com}_{x,y}}
\newcommand{\comVelPlaneLowerBound}{\ubar{\dot{\com}}_{x,y}}
\newcommand{\comVelPlaneUpperBound}{\bar{\dot{\com}}_{x,y}}
\newcommand{\comAcc}{\ddot{\com}}
\newcommand{\comAccPlane}{\ddot{\com}_{x,y}}

\newcommand{\torque}{\bs{\tau}}
\newcommand{\ankleTorque}{{^a\bs{\tau}}}

% The arguments are: origin(or a point), surface normal
\newcommand{\plane}[2]{\Pi(#1, #2)}

% It rotates from frame #1 to #2. 


\newcommand{\eeVel}{\dot{\bs{x}}}

%%%%%%%%%%%%%%%%%%%%%%%%% Momenta
\newcommand{\pGain}{k_p}
\newcommand{\dGain}{k_d}
\newcommand{\iGain}{k_i}

% #1:index of the pole
\newcommand{\pole}[1]{p_{#1}}

% The system momentum matrix
\newcommand{\smm}{
  \textcolor{bodyVelColor}{
    \bs{A}
  }
}

\newcommand{\systemMomenta}{
  \textcolor{bodyVelColor}{
    \bs{h}
  }
}


% Composite rigid body inertia
%%% #1:reference frame, #2:link index
\newcommand{\crbInertia}[2]{
  \textcolor{spatialVelColor}{
    \spatialGInertia[#1][#2]
  }
}
\newcommand{\systemVel}{
  \textcolor{bodyVelColor}{
    \bs{V}
  }
}
\newcommand{\systemJacobian}{
  \textcolor{bodyVelColor}{
    \bs{J}
  }
}

\newcommand{\momentum}{\bs{h}}
\newcommand{\momentumd}{\dot{\bs{h}}}

%%% #1 reference frame, #2 link index
\newcommand{\momentumFrame}[2]{\bs{h}_{#1#2}}

%%% The canonical momenta: 
\newcommand{\canoMomenta}{{\bs{h}_{J}}}

%%% The system momentum per link
\newcommand{\sMomentum}{{\bs{h}_{s}}}

%%% The contact momentum per contact
\newcommand{\cMomentum}{{\bs{h}_{c}}}

% #1:reference frame, #2:contact index
\newcommand{\cMomentumDef}[2]{
  \equivalentMass_{#2} \spatialVel{#1}{#2}
}

\newcommand{\cMomentumDefDef}[2]{
  {\equivalentMassDef{#1}{#2}} \spatialVel{#1}{#2}
}


\newcommand{\cmmMomenta}{
  \textcolor{bodyVelColor}{
  \bs{h}_{\cmmFrame}
 }
}
\newcommand{\cmmInertia}{
  \textcolor{bodyVelColor}{
    \bs{I}_{\cmmFrame}
  }
}
\newcommand{\cmmInertiaDef}{
  \agg{i}{\dof}{\spatialGInertia{\cmmFrame}{i}}
}

\newcommand{\cmmInertiaDefTwo}{
  \agg{i}{\dof}{\spatialGInertia{\cmmFrame}{i}}
}

\newcommand{\cmmInertiaDetail}{
  \textcolor{bodyVelColor}{
    \bs{I}_{\cmmFrame}
  }
}
\newcommand{\angularMomentum}{\mathcal{L}}
\newcommand{\angularMomentumD}{\dot{\angularMomentum}}
\newcommand{\linearMomentum}{P}
\newcommand{\linearMomentumD}{\dot{\linearMomentum}}

\newcommand{\impactDuration}{\delta t}



\newcommand{\nextState}[1]{{#1}_{t_{k+1}}}
\newcommand{\nextStateB}[2]{{#1}_{#2,t_{k+1}}}
\newcommand{\nextStateTwo}[1]{{#1(t_{k+1})}}


\newcommand{\selectionMatrix}{S}

\newcommand{\uncertain}[1]{\widetilde{#1}}

\newcommand{\samplingPeriod}{\Delta t}


% %%%%%%%%%%%%%%%%%%%%%%%%% EQUATION VARIABLES:

\newcommand{\cmm}{
  % \textcolor{bodyVelColor}{
    A_{\cmmFrame}(\jangles)
  % }
}

\newcommand{\cmmLinear}{A_{\bs{v}\cmmFrame}}
\newcommand{\cmmXY}{A^{x,y}_{\cmmFrame}}

\newcommand{\cmmDot}{
  \textcolor{bodyVelColor}{
  \dot{A}_{\cmmFrame}(\jointAngles)
}
}

\newcommand{\cmmPlane}{
  \textcolor{bodyVelColor}{
    A_{\cmmFrame_{x,y}}
  }
}

\newcommand{\jump}{\Delta}
\newcommand{\impactduration}{\delta t}
\newcommand{\sampletime}{\Delta t}
\newcommand{\eeposition}{\bs{x}}
\newcommand{\eevelocity}{\bs{\dot{x}}}
\newcommand{\eeacceleration}{\bs{\ddot{x}}}
\newcommand{\optVars}{\bs{u}}
\newcommand{\argmin}[2]{\underset{#1}{\text{argmin}}{(#2)}}
\newcommand{\argmax}[2]{\underset{#1}{\text{argmax}}{(#2)}}


\newcommand{\gravity}{\mathbf{g}}

\newcommand{\jointState}{\mathbf{q}}
\newcommand{\jointVel}{\dot{\jointState}}
\newcommand{\jointAcc}{\ddot{\jointState}}


\newcommand{\Id}{\mathbb{I}}


% The center of pressure (COP) point
\newcommand{\cop}{{\bs{c}_{p}}}

% The zero-tilting moment point (ZMP)
\newcommand{\zmp}{{\text{\bf z}}}

\newcommand{\height}[1]{{d_{#1}}}

\newcommand{\zmpd}{{\dot{\zmp}}}

\newcommand{\zmpx}{{\zmp_x}}
\newcommand{\zmpxd}{{\zmpd_x}}

\newcommand{\zmpy}{{\zmp_y}}
\newcommand{\zmpyd}{{\zmpd_y}}

% The centroidal moment point/pivot (CMP)
\newcommand{\cmp}{{\bs{c}_{mp}}}

% The divergent component of motion (DCM) or the instantaneous capture point (ICP)
\newcommand{\dcm}{{\bs{\xi}}}
\newcommand{\dcmd}{{\dot{\bs{\xi}}}}
\newcommand{\dcmDef}{{\com + \frac{\comVel}{\pendulumC}}}
\newcommand{\dcmxDef}{{\com_x + \frac{\comVel_x}{\pendulumC}}}

\newcommand{\CoM}{\boldsymbol{c}}
\newcommand{\VRP}{\boldsymbol{v}}
\newcommand{\eCMP}{\boldsymbol{e}}
\newcommand{\DCM}{\boldsymbol{\xi}}
\newcommand{\PluckerPose}{\boldsymbol{p}}
\newcommand{\ZMP}{\mathbf{p}_z}

% The center of mass (COM)
\newcommand{\com}{{\text{\bf c}}}
\newcommand{\comd}{{\dot{\com}}}
\newcommand{\comdxy}{\comd_{xy}}
\newcommand{\comdPlane}{{\dot{\com}_{x,y}}}
\newcommand{\comddPlane}{{\ddot{\com}_{x,y}}}

\newcommand{\picomdxy}{\postImpact{\comd}_{xy}}
\newcommand{\picomdSet}{\set{\postImpact{\comd}_{xy}}}
\newcommand{\picomdVertex}[1]{\set{\postImpact{\comd}_{xy},#1}}
\newcommand{\comdd}{{\ddot{\com}}}

\newcommand{\comx}{{\bs{c}_x}}
\newcommand{\comxd}{{\comd_x}}

\newcommand{\comy}{{\bs{c}_y}}
\newcommand{\comyd}{{\comd_y}}

% stance foot location
\newcommand{\stanceFL}{{\bs{p}_s}}
% capture foot location
\newcommand{\captureFL}{{\bs{p}_c}}

% The arguments are: (1) starting point (2) the ending point. For instance: \vec{a} - \vec{b} = \pointer{b}{a}. 
\newcommand{\pointer}[2]{\overrightarrow{#1#2}}
\newcommand{\pointerDef}[2]{\bs{#2}  - \bs{#1}}

\newcommand{\pointerPerp}[2]{\overrightarrow{#1#2}_{\perp}}


\newcommand{\impulse}{\bs{\iota}} %TODO similar to identityMatrix
\newcommand{\impulseScalar}{\iota} %TODO similar to identityMatrix

\newcommand{\specialJac}{\mathcal{J}}
\newcommand{\specialJacTwo}{\mathcal{C}}
\newcommand{\constant}{C}

\newcommand{\constantVector}{\bs{c}}

\newcommand{\supportPolygon}{\set{\mathcal{S}}}
\newcommand{\zmpArea}{{\set{[\zmp]}}}
\newcommand{\comArea}{{\set{[\com]}}}

\newcommand{\comdArea}{\set{[\comd_{xy}]} }

\newcommand{\controlArea}{{\mathcal{S}_{\com\zmp}}}
\newcommand{\contactArea}{{\mathcal{S}_{\contactPoint}}}

% #1 joint number
\newcommand{\svaJointJac}[1]{\bf{\Phi}_{#1}}
\newcommand{\svaJointJacDef}[1]{\bf{\Phi}_{#1}}


\newcommand{\jangles}{\mathbf{q}}
\newcommand{\jvelocities}{\mathbf{\dot{q}}}
\newcommand{\jaccelerations}{\mathbf{\ddot{q}}}
\newcommand{\jtorques}{\boldsymbol{\tau}}
\newcommand{\omegaVector}{\boldsymbol{\omega}}
\newcommand{\tsInertia}{\boldsymbol{\Lambda}}
\newcommand{\jacobian}{J}
\newcommand{\comJacobian}{J_{\com}}
\newcommand{\comJacobiand}{\dot{J}_{\com}}

\newcommand{\jacobianVector}{\bs{J}}

\newcommand{\jacobiand}{\dot{\bs{J}}}
\newcommand{\Kinv}{\bs{K}^{-1}}
\newcommand{\JEstimate}{\bs{K}_{\jvelocitiesJump}}
\newcommand{\JpinvEstimate}{\Kinv_{\jvelocitiesJump}}
\newcommand{\LambdaINVEstimate}{\bs{K}_{\impulse}}
\newcommand{\LambdaEstimate}{\Kinv_{\impulse}}

\newcommand{\weightMatrix}[1]{\mathbb{#1}}
\newcommand{\weight}{w}

\newcommand{\weightmatrixH}{\mathbf{H}}
\newcommand{\weightmatrixW}{\mathbf{W}}
\newcommand{\weightmatrixQ}{\mathbf{Q}}


\newcommand{\matrixA}{\mathbf{A}}

\newcommand{\quantityJump}{\jump \quantity}
\newcommand{\jtorquesJump}{\jump \jtorques}
\newcommand{\janglesJump}{\jump \jangles}
\newcommand{\jvelocitiesJump}{\jump \jvelocities}
\newcommand{\eevelocityJump}{\jump \eevelocity}
\newcommand{\forceJump}{\jump \force}
\newcommand{\dcmJump}{\jump \dcm}
\newcommand{\zmpJump}{\jump \zmp}

\newcommand{\ubar}[1]{\underaccent{\bar}{#1}}
\newcommand{\jangle}{q}

\newcommand{\lowerBound}[1]{\underline{#1}}
\newcommand{\upperBound}[1]{\overline{#1}}
\newcommand{\doubleBound}[1]{\ubar{\bar{#1}}}




\newcommand{\jvelocitiesUpperBound}{\mathbf{\dot{\bar{\jangle}}}}
\newcommand{\jvelocitiesLowerBound}{\mathbf{\dot{\ubar{\jangle}}}}
\newcommand{\jvelocitiesJumpUpperBound}{\jump \jvelocitiesUpperBound}
\newcommand{\jvelocitiesJumpLowerBound}{\jump \jvelocitiesLowerBound}

%%%%%%%%%%%%%%%%%%%%%%%%%

% \newtheorem{theorem}{Theorem}
% \newtheorem{Lemma}{Lemma}
\newtheorem{remark}{Remark}[section]
\newtheorem{assumption}{assumption}
\newtheorem{hypothesis}{hypothesis}
\newtheorem{proposition}[theorem]{Proposition}
\newtheorem{property}[section]{property}
% \newtheorem{property}[theorem]{property}
% \newtheorem{corollary}[theorem]{corollary}
% \newtheorem{Example}[theorem]{Example}
% \newtheorem{problem}[theorem]{Problem}
% \newtheorem{problem}{Problem}
\newtheorem{Solution}{Solution}
% \newtheorem{Definition}{Definition}


% -----------------------------------------------------------------
\newtheorem{thm}{Theorem}[section]
\newtheorem{cor}[thm]{Corollary}
\newtheorem{lem}[thm]{Lemma}
\newtheorem{prop}[thm]{Proposition}
\theoremstyle{definition}
\newtheorem{defn}[thm]{Definition}
\theoremstyle{remark}
\newtheorem{rem}[thm]{Remark}

%%%%%%%%%%%%%%%%%%%%%%%%%% Def equal

% \usepackage{mathtools}
\newcommand{\defeq}{\vcentcolon=}
\newcommand{\eqdef}{=\vcentcolon}


% Operators
% New operators must defined as such to have them typeset
% correctly. As an example we define the Jacobian:
% -----------------------------------------------------------------
% \DeclareMathOperator{\Jac}{Jac}


% Shortcuts.
% One can define new commands to shorten frequently used
% constructions. As an example, this defines the R and Z used
% for the real and integer numbers.
% -----------------------------------------------------------------
\def\RR{\mathbb{R}}
\def\ZZ{\mathbb{Z}}

\newcommand{\RRv}[1]{\mathbb{R}^{#1}}
\newcommand{\RRm}[2]{\mathbb{R}^{#1 \times #2}}

%%%%%%%%%%%%%%%%%%%%%%%%%%% Impact-handling

\newcommand{\current}[1]{{#1}_{t_{k}}}
\newcommand{\currentTwo}[1]{#1(t_{k})}
\newcommand{\nextTwo}[1]{{#1(t_{k+1})}}
\newcommand{\previous}[1]{{#1}_{t_{k-1}}}


\newcommand{\surfaceNormal}{\bs{n}}
\newcommand{\projection}{P}

\newcommand{\projectionDef}[1]{{#1}{#1}^{\top}}

%%%%%%%%%%%%%%%%%%%%%%%%%%% Algorithm

% \usepackage[ruled,lined,linesnumbered]{algorithm}
% \usepackage[ruled,lined]{algorithm}
% \usepackage[noend]{algpseudocode}
% \usepackage[ruled,vlined,linesnumbered]{algorithm2e}
% \include{pythonlisting}

%%%%%%%%%%%%%%%%%%%%%%%%%%% C++ code

\definecolor{dkgreen}{rgb}{0,0.6,0}
\definecolor{gray}{rgb}{0.5,0.5,0.5}
\definecolor{mauve}{rgb}{0.58,0,0.82}

\lstset{frame=tb,
  language=C++,
  backgroundcolor=\color{black!5}, % set backgroundcolor
  basicstyle=\footnotesize,% basic font setting
  aboveskip=3mm,
  belowskip=3mm,
  showstringspaces=false,
  columns=flexible,
  numbers=none,
  numberstyle=\tiny\color{gray},
  keywordstyle=\color{blue},
  commentstyle=\color{dkgreen},
  stringstyle=\color{mauve},
  breaklines=true,
  breakatwhitespace=true,
  tabsize=3
}


% \NewDocumentCommand{\cpp}{v}{%
% \texttt{\textcolor{blue}{#1}}%
% }

% %%%%%%%%%%%%%%%%%%%%%%%%%%%%% New reference command:
% \NewDocumentCommand{\probRef}{v}{%
% Problem~\ref{#1}}

% \NewDocumentCommand{\secRef}{v}{%
% Sec.~\ref{#1}}

% \NewDocumentCommand{\remarkRef}{v}{%
% Remark~\ref{#1}}

% \NewDocumentCommand{\tableRef}{v}{%
% Table~\ref{#1}}

% \NewDocumentCommand{\algRef}{v}{%
% Algorithm~\ref{#1}}

% \NewDocumentCommand{\lemmaRef}{v}{%
% Lemma~\ref{#1}}

% \NewDocumentCommand{\theoremRef}{v}{%
% Theorem~\ref{#1}}

% \NewDocumentCommand{\figRef}{v}{%
%   Fig.~\ref{#1}}

% \NewDocumentCommand{\appRef}{v}{%
%   Appendix~\ref{#1}}

% \NewDocumentCommand{\exRef}{v}{%
% Experiment.~\ref{#1}}

% \NewDocumentCommand{\simRef}{v}{%
% Simulation.~\ref{#1}}


% \NewDocumentCommand{\orderedS}\TwoArgs{mm}{v}{%
%   <#1, #2, #3>
% }
% \NewDocumentCommand\orderedS{m{}m{}O{}}{
% % <#1, #2, #3>
% }
% \NewDocumentCommand\TwoOptOfThree{O{}O{}m}%
%   {Text with #3 and perhaps #1 and #2}
% \NewDocumentCommand\orderedTwoS{mm}%
%   {$<$#1,#2$>$}
% \NewDocumentCommand\orderedThreeS{mmm}%
%   {$<$#1,#2,#3$>$}

% \usepackage{xcolor}
\definecolor{svaColor}{RGB}{7, 160, 2}
\definecolor{uncertainColor}{RGB}{255, 0, 0}

\definecolor{spatialVelColor}{RGB}{189, 38, 96}
\definecolor{geometricColor}{RGB}{66, 126, 245}
\definecolor{bodyVelColor}{RGB}{13, 191, 191}

\definecolor{velColor}{RGB}{50, 205, 50}


\colorlet{svaColorVariable}{black}%{svaColor}
\colorlet{geometricColorVariable}{geometricColor}

% %%%%%%%%%%%%%%%%%%%%%%%%% comment
\newif\ifBlockComment
\BlockCommentfalse %\BlockCommenttrue
% %%%%%%%%%%%%%%%%%%%%%%%%% 
